\documentclass[10pt,a4paper]{amsart}
\usepackage[margin=19mm]{geometry}
\usepackage{amsmath,amssymb,mathtools}
\usepackage[scr=rsfs,cal=euler]{mathalfa}
\usepackage{xfrac}
\usepackage[unicode,hidelinks,bookmarks=false]{hyperref}
\newcommand{\ie}{i.e.\ }
\newcommand{\cf}{cf.\ }
\newcommand{\resp}{respectively\ }
\newcommand{\Subsection}{\subsection}
\providecommand{\nobreakdash}{\nobreak\hbox{-}}
\setlength{\emergencystretch}{2em}
\multlinegap0pt
\def\AC{\mathcal{A}}
\def\IC{\mathcal{I}}
\def\NC{\mathcal{N}}
\def\PC{\mathcal{P}}
\def\MC{\mathcal{M}}
\def\BC{\mathcal{B}}
\def\DC{\mathcal{D}}
\def\LC{\mathcal{L}}
\def\TC{\mathcal{T}}
\def\XC{\mathcal{X}}
\def\CC{\mathcal{C}}
\def\UC{\mathcal{U}}
\def\EC{\mathcal{E}}
\def\FC{\mathcal{F}}
\def\HC{\mathcal{H}}
\def\GC{\mathcal{G}}
\def\C{\mathbf{C}}
\def\E{\mathbf{E}}
\def\N{\mathbf{N}}
\def\P{\mathbf{P}}
\def\Q{\mathbf{Q}}
\def\R{\mathbf{R}}
\def\Z{\mathbf{Z}}
\def\x{\mathbf{x}}
\def\y{\mathbf{y}}
\def\1{\mathbf{1}}
\def\z{\mathbf{z}}
\def\w{\mathbf{w}}
\def\u{\mathbf{u}}
\def\v{\mathbf{v}}
\def\n{\mathbf{n}}
\def\e{\mathbf{e}}
\def\tr{\rm{tr}}
\def\supp{\rm{supp}}
\def\sym{\rm{sym}}
\def\sgn{\rm{sgn}}
\def\Int{\rm{Int}}
\def\div{\rm{div}}
\def\al{\alpha}
\def\be{\beta}
\def\pa{\partial}
\def\ep{\epsilon}
\def\de{\delta}
\def\ga{\gamma}
\def\ka{\varkappa}
\newtheorem{prop}{Proposition}[section]
\newtheorem{theorem}{Theorem}[section]
\newtheorem{lemma}{Lemma}[section]
\newtheorem{exer}{Exercise}[section]
\newtheorem{corollary}{Corollary}
\newtheorem{remark}{Remark}
\newcommand{\la}{\lambda}
\newcommand{\si}{\sigma}
\newcommand{\ph}{\varphi}
\newcommand{\om}{\omega}
\newcommand{\Ga}{\Gamma}
\newcommand{\De}{\Delta}
\newcommand{\Th}{\Theta}
\newcommand{\Si}{\Sigma}
\newcommand{\Om}{\Omega}
\newcommand{\La}{\Lambda}
\newcommand{\lan}{\langle}
\newcommand{\ran}{\rangle}
\begin{document}
\title[Quantum filtering and propagation of chaos for open quantum ]{Quantum filtering and propagation of chaos for open quantum systems, with applications to quantum feedback control and quantum mean-field games}
\author{Moscow Center of Fundamental and Applied Mathematics, Leninskie Gory, 1, 119234 Moscow, Russia; Vassili N. Kolokoltsov}
\date{}
\maketitle
\noindent\emph{Source and licence.} Quantum filtering and propagation of chaos for open quantum systems, with applications to quantum feedback control and quantum mean-field games. Version arXiv:2607.08507v1 (CC BY 4.0). This material is distributed under the Creative Commons Attribution 4.0 International licence, http://creativecommons.org/licenses/by/4.0/, as recorded in the arXiv OAI licence metadata for this exact version and shown on the abstract page https://arxiv.org/abs/2607.08507v1. Full rights notice: licenses/arxiv-2607.08507-CC-BY-4.0.txt.\par\smallskip
\noindent\textit{Editorial scope.} Counted unit: the original \texttt{theorem} environment \texttt{thBeleqcount} at source lines 2724--2747 together with its complete proof at lines 2749--2757; the delivered document keeps source lines 372--3015 so that every referenced label and every citation resolves inside it. Native editable source is the paper's own concatenated TeX; the AMS wrapper is editorial formatting only. No publisher class, upstream script or unrelated artwork is redistributed.
\begin{equation}
\label{Lindstoch}
d\ga(t)=-i[H,\ga(t)] \, dt +\LC_L \ga(t) \, dt +(L\ga(t)+\ga(t) L^*) dY(t),
\end{equation}
with
\[
\LC_L\ga =L\ga L^*-\frac12 L^*L\ga -\frac12 \ga L^*L
=L\ga L^*-\frac12 \{L^*L,\ga\}.
\]

We shall mostly use everywhere such concise form, only occasionally writing the full version
(when it seems to be necessary for clarity).
For finite-dimensional $\HC$, this is a simple linear matrix-valued equation, so that
the well-posedness of the corresponding Cauchy problem follows from the standard results of Ito's calculus.
For general $\HC$ the situation is not so obvious even for bounded $H,L$ and will be discussed later on.

Using Ito's rule one directly derives a stochastic equation for the normalised state
$\rho=\ga/{\tr} \, \ga$:
\begin{equation}
\label{Lindstochnorm1}
d\rho(t)=-i[H,\rho(t)]\, dt+\LC_L \rho (t)\, dt
+[L\rho(t)+\rho(t) L^*-\rho(t)\, {\tr} \, (L\rho(t)+\rho(t) L^*) ] dB(t),
\end{equation}
where
\begin{equation}
\label{eqdefinnov}
dB(t)=dY(t)-2{\tr} \, (\rho(t), L_S) \, dt
\end{equation}
is the Ito process, referred to as the {\it innovation process}.

According to Girsanov's theorem, by an equivalent change of measure one can turn $B(t)$
to the standard Brownian motion (BM). And in fact, one usually chooses $Y$ to be a BM when dealing with a linear
equation and $B$ when dealing with the normalised one.

The evolutions described by \eqref{Lindstochnorm1} and  \eqref{Lindstoch},
their various versions and extensions (including counting evolutions \eqref{eqBeleqcountj}),
are the main objects of studies in this paper.
A derivation of these equations from the point of view of continuous quantum measurement
will be given in Section \ref{secder}.

\begin{remark}
As already mentioned above,
ignoring the noise term in both \eqref{Lindstochnorm1} and \eqref{Lindstoch}
(setting $Y=0$ or $B=0$) leads to
the famous Lindblad equations describing general quantum dynamic semigroups. Therefore, some physicists suggested
equations \eqref{Lindstochnorm1} and \eqref{Lindstoch} as a way of stochastic unravelling of quantum dynamic semigroups,
see numerous references in \cite{Mora13} or \cite{Breuer}.
\end{remark}

Let us make some comments about the evolution of traces.
Equation \eqref{Lindstochnorm1} is best suited for dealing with matrices $\ga$ with unit trace.
A convenient extension to arbitrary traces can be written as
\begin{equation}
\label{Lindstochnorm1m}
d\rho(t)=-i[H,\rho(t)]\, dt +\LC_L \rho(t)\, dt
+[\rho(t)L^*+L\rho(t)-\frac{\rho(t)}{{\tr} \, \rho(t)} \, {\tr} (\rho(t)(L^*+L))]\, dB(t).
\end{equation}
It is seen directly from this equation that $d\, {\tr} \, \rho(t)=0$ and hence the trace is preserved
by evolution \eqref{Lindstochnorm1m} (as long as it is well-posed of course).
For equation \eqref{Lindstochnorm1} one gets
 \begin{equation}
\label{Lindstochnorm1tr}
d\, {\tr} \, \ga(t)= {\tr} (\ga(t)(L^*+L))(1-  {\tr} \, \ga(t)) \, dB(t).
\end{equation}
Thus  evolution \eqref{Lindstochnorm1} does not preserve traces in general. But
${\tr} \, \ga(t)=1$ is a solution of \eqref{Lindstochnorm1tr}. Hence, if the initial
condition has  ${\tr} \, \ga_0=1$ and the equation \eqref{Lindstochnorm1} is well posed, then
their solutions do preserve the trace.

\begin{remark}
In this paper we reduce attention to finite-dimensional noises $B(t)$ or $Y(t)$.
In literature, see e.g. \cite{MoraRebo}, one can find also the corresponding equations
with infinite-dimensional noises. The author believes that most of the results presented
here can be naturally extended to infinite $n$, if carefully looking at the convergence
of the related infinite sums. We decided not to overload this paper
with related technicalities.
\end{remark}

A remarkable property of evolutions \eqref{Lindstochnorm1} and  \eqref{Lindstoch} is the fact that
they preserve pure states. Namely, direct application of Ito's lemma shows that if vectors
 $\chi(t)\in \HC$ evolve according to the {\it linear Belavkin quantum filtering equation
 for non-normalized pure states}
\begin{equation}
\label{eqqufiBlins}
d\chi(t) =-[iH\chi(t) +\frac12 L^*L \chi(t) ]\,dt+L\chi(t) \, dY(t),
\end{equation}
then the decomposable density matrices $\ga(t)=\chi(t)\otimes \bar \chi(t)$
satisfy equation \eqref{Lindstoch}. In fact, from \eqref{eqqufiBlins} it follows that
\[
d\overline{\chi(t)} =[i \overline{H\chi(t)} -\frac12 \overline{L^*L \chi(t)} ]\,dt+\overline{L\chi(t)}\, dY(t).
\]
Consequently, by Ito's product rule,
\[
d\ga(t)=[-iH\chi(t)\otimes \bar \chi -\frac12 L^*L \chi(t)\otimes \bar \chi (t)]\,dt+L\chi(t)\otimes \bar \chi(t) \, dY(t)
\]
\[
+[i \chi(t)\otimes \overline{H\chi(t)} -\frac12 \chi(t) \otimes \overline{L^*L \chi(t)} ]\,dt
+\chi(t)\otimes \overline{L\chi(t)}\, dY(t)+L\chi(t) \otimes \overline{L\chi(t)} \, dt
\]
\[
=(-iH\ga+i\ga H)\, dt -\frac12 \{L^*L,\ga\} \, dt +(L\ga+\ga L^*) \, dY(t)+L\ga L^* \, dt,
\]
as claimed.

On the other hand, as again follows from Ito's formula,
if $\chi(t)$ satisfy \eqref{eqqufiBlins}
then the corresponding normalised vectors $\phi(t)=\chi(t)/\|\chi(t)\|$ satisfy
the {\it nonlinear Belavkin quantum filtering equation for normalized pure states}
\[
d\phi(t)=(L-(\phi(t), L_S \phi(t)))\phi(t) \, dB(t)
\]
\begin{equation}
\label{eqqufiBnonlins}
-[i(H-(\phi(t), L_S \phi(t)) L_A)
+\frac12 (L-(\phi(t), L_S \phi(t)))^*(L-(\phi(t), L_S \phi(t)))]\phi(t) \, dt,
\end{equation}
and the corresponding decomposable normalised
matrices $\rho(t) =\phi(t)\otimes \bar \phi(t)$ satisfy \eqref{Lindstochnorm1},
with $B$ given by \eqref{eqdefinnov}:
\[
dB(t)=dY(t)-2{\tr} \, (\rho(t), L_S) \, dt=dY(t)-2(\phi(t), L_S \phi(t))\, dt.
\]

\begin{remark}
To see the pure mathematical place of equation \eqref{eqqufiBlins} in the theory of SDEs, we can note that
this is a general linear equation in a Hilbert space preserving the expectation of norm squared,
or where the norm squared of a solution is a local martingale. In fact, writing a linear SDE in the form
$d\chi=A\chi \,dt +L\chi \, dY$ with a given operator $L$ and requiring that $d \, \E \|\chi\| ^2(t)=0$,
we get by Ito's formula (at least, its formal application) that $A$ has a structure given in \eqref{eqqufiBlins}.
\end{remark}

\begin{remark}
\label{remtimedep}
Everything given below have a straightforward extension to the case of time dependent families $L(t)$, as
long as these families are continuous (at least in the strong topology)
 and uniformly bounded. In fact, time dependent families arise necessarily
when reducing the case of unbounded $H$ to the case of bounded (in fact, vanishing) $H$ via the interaction
 representation, see below. If $L$ is time-dependent, all estimates below are valid with $\|L\|=\max_t\|L(t)\|$.
Similarly the theory extends to time-dependent Hamiltonian $H_t$ whenever such family generates a well defined
unitary propagator.
\end{remark}


\begin{remark} The square norm $\|\chi\|^2$ of solutions to \eqref{eqqufiBlins} has physical meaning
analogous to the square norm of the solutions to the standard Schr\"odinger equation:
it describes the probability density of observing corresponding values of the output process $Y(t)$,
see detailed discussion in  \cite{BarchBook}.
\end{remark}


As for the case of equations on density matrices,
it is often convenient
to include equations \eqref{eqqufiBnonlins}
in a more general class of norm preserving evolutions, the simplest version being
\[
d\phi(t)=-[i(H-\langle L_S \rangle_{\phi(t)} L_A)
+\frac12 (L-\langle L_S \rangle_{\phi(t)})^*(L-\langle L_S \rangle_{\phi(t)})]\phi(t) \, dt
\]
\begin{equation}
\label{eqqufiBnonlinsn}
+ (L-\langle L_S \rangle_{\phi(t)})\phi(t) \, dB(t),
\end{equation}
where we introduced the (rather standard) notation for the value
of an operator $A$ in a pure state $\phi$:
\[
\langle A \rangle_{\phi}=\frac{(\phi, A \phi)}{(\phi, \phi)}.
\]
Clearly for $\phi(t)$ of unit norm solutions to equations \eqref{eqqufiBnonlins}
and \eqref{eqqufiBnonlinsn} coincide, but \eqref{eqqufiBnonlinsn} is explicitly norm-preserving
for arbitrary $\phi(t)$, which is not the case for equation \eqref{eqqufiBnonlins}.

It is also insightful to write down equation for $\chi$ in terms of the innovation process $B$:
\begin{equation}
\label{eqqufiBlinsB}
d\chi(t) =-[iH\chi(t) +\frac12 L^*L \chi(t) ]\,dt
+L \chi(t) \, (dB(t)+\langle L+L^*\rangle_{\chi(t)} \, dt).
\end{equation}

In the most important case of a self-adjoint $L$ equations
\eqref{eqqufiBlins} and \eqref{eqqufiBnonlinsn} simplify to the equations
\begin{equation}
\label{eqqufiBlinss}
d\chi(t) =-[iH\chi(t) +\frac12 L^2 \chi(t) ]\,dt+L\chi(t) dY(t),
\end{equation}
and, respectively,
\begin{equation}
\label{eqqufiBnonlinss}
d\phi(t)=-[iH+\frac12 (L-(\phi(t), L \phi(t)))^2]\phi(t) \, dt
+ (L-(\phi(t), L \phi(t)))\phi(t) \, dB(t).
\end{equation}

For working with control problems, it is often convenient to work with the Stratonovich rather
than Ito's integral. Using the standard rule for transforming Ito's differential to Stratonovich,
\[
\si(X(t)) \, dB(t)=\si(X(t)) \circ dB(t)-\frac12 \frac{\pa \si}{\pa X}(X(t)) \si(X(t)) \, dt,
\]
equation \eqref{eqqufiBnonlinss} rewrites in terms of the Stratonovich differential $\circ \, dB$ as
\[
d\phi(t)=-[iH+(L-(\phi(t), L \phi(t)))^2 + (\phi(t), L^2 \phi(t))-(\phi(t), L \phi(t))^2 ]\phi(t) \, dt
\]
\begin{equation}
\label{eqqufiBnonlinssStr}
+ (L-(\phi(t), L \phi(t)))\phi(t) \, \circ \, dB(t).
\end{equation}

Notice finally that equations \eqref{Lindstoch} and \eqref{Lindstochnorm1} coincide when $L$ is anti-Hermitian,
that is $L^*=-L$. However, this case is considered as much less relevant from
physical point of view (see discussion in \cite{BarchBook}).

\subsection{Belavkin's quantum filtering equations: counting case}

For the case of {\it counting observation} the {\it Belavkin quantum filtering equations
for normalised density matrices} have the form of the jump type SDE
\begin{equation}
\label{eqBeleqcountj}
d\rho(t)=(- i[H, \rho(t)] -\frac12 \sum_{j=1}^n \{L_j^*L_j,\rho(t)\}+ {\tr} (L_j\rho(t) L_j^*) \rho(t) )\, dt
+\left(\frac{L_j\rho(t) L_j^*}{{\tr} (L_j\rho(t) L_j^*)}-\rho(t)\right) dN_j(t),
\end{equation}
with the counting process $N_j(t)$ having position dependent intensity ${\tr} (L_j^*L_j\rho)$,
so that the compensated processes $M_j(t)=N_j(t)-\int_0^t {\tr} (L_j^*L_j\rho_s) \, ds$ are martingales.

As usual, we shall write these equations in a shorter form
(also omitting for brevity the explicit dependence of unknown functions on $t$):
 \begin{equation}
\label{eqBeleqcount}
d\rho=(- i[H, \rho] -\frac12 \{L^*L,\rho\}+ {\tr} (L\rho L^*) \rho ) \, dt
+\left(\frac{L\rho L^*}{{\tr} (L\rho L^*)}-\rho\right) dN(t).
\end{equation}

In terms of the compensated processes $M_j(t)$, this equation rewrites as
\begin{equation}
\label{eqBeleqcountm1}
d\rho=- i[H, \rho]\, dt +(L\rho L^*-\frac12 \{L^*L,\rho\} ) \, dt
+\left(\frac{L\rho L^*}{{\tr} (L\rho L^*)}-\rho\right) dM(t).
\end{equation}

As in the diffusive case, equation \eqref{eqBeleqcount} is (at least formally) trace preserving
only for operators of unit trace. It is often convenient to extend it to a general class of equations
that is trace preserving for all traces:

\begin{equation}
\label{eqBeleqcounttp}
d\rho=(- i[H, \rho] -\frac12 \{L^*L,\rho\}+ {\tr} (L\rho L^*) \frac{\rho}{{\tr}\, \rho} ) \, dt
+\left( L\rho L^*  \frac{{\tr}\, \rho}{{\tr} (L\rho L^*)}-\rho\right) dN(t).
\end{equation}
Equations \eqref{eqBeleqcount} and \eqref{eqBeleqcounttp} coincide for ${\tr}\,\rho=1$,
but \eqref{eqBeleqcounttp} implies $d\,{\tr}\, \rho=0$ for all trace-class $\rho$.


In equation \eqref{eqBeleqcount} (written in the most common form in the literature)
the driving noise $N(t)$ is itself
position dependent, which makes it not quite standard.
Moreover, a left continuous version of an integrand is not stressed there.
A convenient rigorous
form can be written as follows. Let $K$ be any constant such that $ \| L_j^*L_j \| <K$ for all $j$.
Let $N(dx \, dt)=(N_1(dx \, dt), \cdots, N_n(dx\, dt))$ be a collection of standard independent
Poisson random measures (or corresponding compound Poisson processes)
on $[0,K]\times \R_+$ (with Lebesgue measure
as intensity on $\R_+$). A mathematically rigorous form of \eqref{eqBeleqcount} is as follows:
\[
\rho (t)=\int_0^t (- i[H, \rho(s)] -\frac12 \{L^*L,\rho(s)\}+ {\tr} (L\rho(s) L^*) \rho(s) )\, ds
\]
\begin{equation}
\label{eqBeleqcountind}
+\int_0^t \int_0^K \left(\frac{L\rho(s-) L^*}{{\tr} (L\rho(s-) L^*)}-\rho(s-)\right)
\1(x< {\tr} (L^*L\rho(s-))) N(dx\, ds),
\end{equation}
or, in full version,
 \[
\rho (t)=\int_0^t (- i[H, \rho(s)] -\frac12 \sum_j \{L_j^*L_j,\rho(s)\}
+ \sum_j {\tr} (L_j\rho(s) L_j^*) \rho(s) )\, ds
\]
\[
+\int_0^t \sum_j \int_0^K \left(\frac{L_j\rho(s-) L_j^*}{{\tr} (L_j\rho(s-) L_j^*)}-\rho(s-)\right)
\1(x< {\tr} (L_j^*L_j\rho(s-))) N_j(dx\, ds).
\]

The counting process $N(t)$ from \eqref{eqBeleqcount} is recovered by the formula
\[
N(t)=\int_0^t \int_0^K \1(x< {\tr} (L^*L\rho(s-))) N(dx \, ds).
\]
Talking about equation \eqref{eqBeleqcount}, we shall always mean its rigorous version \eqref{eqBeleqcountind}.

%Yet another equivalent and insightful form of equation \eqref{eqBeleqcountind} is as follows.
%Let $N^K(t)=(N^K_1(t), \cdots, N^K_n(t))$
%be a collection of independent standard Poisson processes of intensity $K$ and
%$\xi_1, \xi_2, \cdots $ a sequence of independent random variable
%distributed uniformly on $[0,K]$.  Then \eqref{eqBeleqcountind} can be written as
%\[
%\rho (t)=\int_0^t (- i[H, \rho(s)] -\frac12 \{L^*L,\rho(s)\}+ {\tr} (L\rho(s) L^*) \rho(s) )\, ds
%\]
%\begin{equation}
%\label{eqBeleqcountind1}
%+\sum_{j=1}^{N_t^K} \left(\frac{L\rho(\tau_j-) L^*}{{\tr} (L\rho(\tau_j-) L^*)}-\rho(\tau_j-)\right)
%\1(\xi_j< {\tr} (L^*L\rho(\tau_j-))),
%\end{equation}
%where $\tau_j$ are the times of jumps of the process $N_t^K$.

As in the case of diffusive measurements, dynamics \eqref{eqBeleqcountind} preserves the set of pure states.
Namely, if $\phi(t)$ satisfy the equation
\begin{equation}
\label{eqcountpure}
d \phi(t)=-\left(iH +\frac12 (L^*L-\|L\phi(t)\|^2)\right)\phi(t) \, dt
+\left(\frac{L\phi(t)}{\|L\phi(t)\|}-\phi(t)\right)  dN(t),
\end{equation}
then $\rho_t=\phi_t\otimes \bar \phi_t$ satisfies equation \eqref{eqBeleqcountm1}.
In fact, $\|L\phi\|^2={\tr}\, (L\rho L^*)$, and,
 by Ito's product rule $dN(t)dN(t)=dN(t)$, it follows that (omitting argument $t$ for brevity)
\[
d\rho=-i[H,\rho] \, dt
-\frac12 (L^*L-\|L\phi\|^2)\phi \otimes \bar \phi\, dt
- \frac12 \phi\otimes  (\overline{L^*L}-\|L\phi\|^2)\bar \phi\, dt
\]
\[
+\left(\frac{L\phi}{\|L\phi\|}-\phi\right) \otimes \bar \phi \, dN(t)
+\phi \otimes \left(\frac{\overline{L\phi}}{\|L\phi\|}-\bar \phi\right) \, dN(t)
\]
\[
+\left(\frac{L\phi}{\|L\phi\|}-\phi\right)
\times \left(\frac{\overline{L\phi}}{\|L\phi\|}-\bar \phi\right)\, dN(t)
\]
\[
=-i[H,\rho] \, dt-\frac12 \{L^*L,\rho\} \, dt
+{\tr}\, (L\rho L^*) \rho  \, dt
+ \left(\frac{L\rho L^*}{{\tr} (L\rho L^*)}-\rho\right) dN(t),
\]
as claimed.

\begin{remark}For readers not well acquainted with stochastic calculus of jump process,
 let us give some intuition for the product rule $dN(t)\,dN(t)=dN(t)$ for Poisson processes
 with any (even variable) intensity. If processes $X_i(t)$, $i=1,2$, are defined by
 their differentials as $dX_i(t)=(F_i(X_i(t-))-X_i(t-)) dN(t)$, this means that at times $\tau_j$
 of the jumps of $N(t)$ they jump from $X_i(\tau_j-)$ to
 \[
 F_i(X_i(\tau_j-))= X_i(\tau_j-)+[F_i(X_i(\tau_j-))-X_i(\tau_j-)].
 \]
 Then $X_1(t)X_2(t)$ at these moments jumps from $X_1(\tau_j-)X_2(\tau_j-)$
 to $F_1(X_1(\tau_j-))F_2(X_2(\tau_j-))$, and thus this process has to have the differential
 \begin{equation}
 \label{eqprodjump}
 d(X_1(t)X_2(t))=\left( F_1(X_1(\tau_j-))F_2(X_2(\tau_j-))-X_1(\tau_j-)X_2(\tau_j-)\right) dN(t).
 \end{equation}
 And that is what we get from Ito's product rule. In fact, it implies that
 \[
 dX_1(t)dX_2(t)=X_1(t) \, dX_2(t)+dX_2(t) \, dX_1(t)+dX_1(t)\, dX_2(t),
 \]
 which equals, by the rule $dN(t)\,dN(t)=dN(t)$, to
 \[
\bigl[ X_1(t)(F_2(X_2(t-))-X_2(t-))+X_2(t) (F_1(X_1(t-))-X_1(t-))
 \]
 \[
 +(F_1(X_1(t-))-X_1(t-)) (F_2(X_2(t-))-X_2(t-))\bigr] dN(t),
 \]
 that is, exactly \eqref{eqprodjump}.
\end{remark}

Again analogously to the diffusive case, equation \eqref{eqBeleqcount} can be obtained
by normalization from an apparently linear equation. Namely, one checks by a straightforward
application of Ito's formula that if $\ga(t)$ satisfy the equation (omitting argument $t$ for brevity)
 \begin{equation}
\label{eqBeleqcountlin}
d\ga=(- i[H, \ga] -\frac12 \{L^*L,\ga\}+ L\ga L^* ) dt
+(L\ga L^*-\ga) (dN(t)-dt),
\end{equation}
then ${\tr \, \ga}$ satisfies the equations
\[
d\, {\tr}\, \ga={\tr} (L\ga L^*-\ga) \, (dN(t)-dt),
\]
or
\[
d\, {\tr}\, \ga=\sum_j {\tr} (L_j\ga L_j^*-\ga) \, (dN_j(t)-dt)
\]
in the full version, and
\[
d\frac{1}{{\tr}\, \ga}=\left(\frac{1}{{\tr}\,(L\ga L^*)}-\frac{1}{{\tr}\, \ga }\right)
\left(dN(t)-\frac{{\tr}\, (L\ga L^*)}{{\tr}\, \ga} dt\right),
\]
implying that $\ga/{\tr}\,\ga$ satisfies \eqref{eqBeleqcount}.

\begin{remark}
Unlike diffusive case, equation \eqref{eqBeleqcountlin} is not exactly linear,
as the intensity of $N(t)$ depends on $\ga$.
\end{remark}

Again pure states are preserved under \eqref{eqBeleqcountlin}.
In fact, if vectors $\chi(t)$ in $\HC$ satisfy the equation
\begin{equation}
\label{eqcountpurelinpu}
d \chi(t)=-\left(iH +\frac12 (L^*L-1)\right)\chi(t) \, dt
+(L-1)\chi(t)  dN(t),
\end{equation}
then $\ga(t)=\chi(t)\otimes \bar \chi(t)$ satisfy \eqref{eqBeleqcountlin}
 and $\phi(t)=\chi(t)/\|\chi(t)\|$ satisfy  \eqref{eqcountpure}.

Let us note finally that in the important particular case when operator $L$ is unitary,
equation \eqref{eqBeleqcount} becomes linear:

 \begin{equation}
\label{eqBeleqcountunit}
d\rho=- i[H, \rho]  \, dt
+\left(L\rho L^*-\rho\right) dN(t),
\end{equation}
with the standard Poisson process $N(t)$ of unit intensity, in which case it coincides, of course,
with \eqref{eqBeleqcountlin} and writes down in terms of the compensated martingale $M(t)=N(t)-t$ as
 \begin{equation}
\label{eqBeleqcountunit1}
d\rho=(- i[H, \rho] + L\rho L^*-\rho ) \, dt
+\left(L\rho L^*-\rho\right) dM(t).
\end{equation}

\subsection{SDEs and PDEs}

Apart from SDEs, the standard alternative way of describing Markov processes
is via PDEs (partial differential equations) or sometimes $\Psi$DEs (pseudo-differential equations).
  These two descriptions are linked via Ito's
formula. Namely, if $X(t)$ is a (time homogeneous) Markov process in some metric space $S$,
starting with $x=X(0)$ and solving
certain SDE, one defines the corresponding Markov transition operators  $T_tf(x)=\E_x f(X(t))$ on an
appropriate space of bounded functions $f$ ($\E_x$ is used to denote the
expectation with respect to the process started in $x$).
When $T_t$ preserves the space of continuous
functions (as will be the case for all examples of this paper), the process and the semigroup
are usually referred to as $C$-Feller. In this case the operators $T_t$ form a semigroup of positive
contractions in the space $C(S)$. By a generator
of the process $X(t)$ one understands the operator $Af=\lim_{t\to 0} (T_tf-f)/t$ (defined on the set of functions,
where this limit exists in some sense). It turns out that, under appropriate regularity assumptions,
the functions $T_tf(x)$ solve the PDE $\pa f/\pa t=Af$. If $A$ is a second order differential operator,
it is referred to as a diffusion operator and the corresponding process is called a diffusion.

For handy future references, let us write now the formal expressions for the generators of all processes
defined by the basic quantum filtering SDEs. They are direct consequences of the corresponding SDEs
and of the formal application of Ito's formula. Of course, only the well-posedness theory for all these equations,
developed further, will make these derivations rigorous.

(1) If $T_tf(\ga)=\E_{\ga} f(\Ga(t))$, where $\Ga(t)$ yields the solution
to the Cauchy problem of SDE \eqref{Lindstoch} with the initial condition $\ga$,
then (for sufficiently smooth functions $f$ on the set of density operators)
the function $T_tf(\ga)=f(t,\ga)$ solves the Cauchy problem
\begin{equation}
\label{eqCauchylinfilt}
\frac{\pa f}{\pa t}=\AC_{lin} f, \quad f(0,\ga)=f(\ga),
\end{equation}
where
\begin{equation}
\label{eqgenlinfilt1}
\AC_{lin} f(\ga)= (f'(\ga), -i[H,\ga] +\LC_L \ga) + \sum_j \frac12 (L_j\ga+\ga L_j^*, f''(\ga) (L_j\ga+\ga L_j^*)).
\end{equation}
or in our usual concise notation
\begin{equation}
\label{eqgenlinfilt}
\AC_{lin} f(\ga)= (f'(\ga), -i[H,\ga] +\LC_L \ga) + \frac12 (L\ga+\ga L^*, f''(\ga) (L\ga+\ga L^*)).
\end{equation}
Therefore, this $\AC_{lin}$ is the generator of the diffusion $\Ga(t)$.

(2) Solutions to SDEs  \eqref{Lindstochnorm1} specify the Markov process on the convex
set of density operators  $S(\HC)$, with the generator
being the diffusion operator
\[
\AC f(\rho)=(f'(\rho),-i[H, \rho] -\frac12 \{L^*L,\rho\}+ L\rho L^*)
\]
\begin{equation}
\label{eqdifgener}
\frac12 \bigl(\rho L^* + L\rho-{\tr}\,(\rho L^* +L\rho)  \rho,
f''(\rho) (\rho L^* + L\rho-{\tr}\,(\rho L^* +L\rho)  \rho)\bigr).
\end{equation}

(3) Solutions to \eqref{eqBeleqcountlin} specify the Markov process with the generator
\[
\AC_{count,lin}f(\rho)=-(f'(\rho), i[H, \rho]
+\sum_{j=1}^n \left(\frac12 \{L_j^*L_j,\rho\}+ \rho\right)
\]
 \begin{equation}
\label{eqjumpgenerlin}
+ \sum_{j=1}^n {\tr}\, (L_j^*L_j\rho)\left[f(L_j\rho L_j^*)-f(\rho)\right].
\end{equation}

(4) Solutions to SDE \eqref{eqBeleqcountind}
specify a Markov process with the generator
\[
\AC_{count}f(\rho)=-(f'(\rho), i[H, \rho] +\sum_{j=1}^n (\frac12 \{L_j^*L_j,\rho\}-\rho \, {\tr}\, (L_j^*L_j\rho))
\]
\begin{equation}
\label{eqjumpgenerj}
+\sum_j {\tr}\, (L_j^*L_j\rho) \left[f\left(\frac{L_j\rho L_j^*}{{\tr}\, (L_j^*L_j\rho)}\right)-f(\rho)\right],
\end{equation}
or, in concise notations,
 \[
\AC_{count}f(\rho)=-(f'(\rho), i[H, \rho] + \frac12 \{L^*L,\rho\}-\rho \, {\tr}\, (L^*L\rho))
\]
 \begin{equation}
\label{eqjumpgener}
+ {\tr}\, (L^*L\rho)\left[f\left(\frac{L\rho L^*}{{\tr}\, (L^*L\rho)}\right)-f(\rho)\right].
\end{equation}

\subsection{Quantum filtering for mixed channels}

Though one considers usually either counting or diffusive observations, and we shall do so most of the time,
it makes sense to organise mixed channels of observations that include both type of observations simultaneously.
Looking at operators \eqref{eqdifgener} and \eqref{eqjumpgenerj} it is easy to guess the expression for the generator
that combines both type of observation:

\[
\AC_{mix}f(\rho)
=\sum_{j=k}^{m-1} \, {\tr} \, (L_j\rho L_j^*)  \left[f(\frac{ L_j \rho  L_j^*}{ {\tr} \, (L_j\rho L_j^*)})-f(\rho)\right]
\]
\[
+\left(f'(\rho), -i[H, \rho] -\frac12 \sum_{j=0}^{m-1}\{ L_j^* L_j,\rho\}
+\sum_{j=0}^{k-1} L_j \rho  L_j^*
+\sum_{j=k}^{m-1}  \, {\tr} \, (L_j\rho L_j^*) \rho \right)
\]
\begin{equation}
\label{eqmixgener}
+\frac12 \sum_{j=0}^{k-1}  \left(\rho  L_j^* +  L_j\rho- {\tr} (\rho L_j^* +L_j \rho)\rho,
f''(\rho)  (\rho  L_j^* +  L_j\rho- {\tr} (\rho L_j^* +L_j \rho) \rho)\right).
\end{equation}

In this paper we reduce our attention to simplified mixed channels,
when all $L_j$ with $j\ge k$ are unitary,
in which case $\AC_{mix}$ becomes
\[
\AC_{mix}f(\rho)
=\sum_{j=k}^{m-1} \left[f(L_j \rho  L_j^*)-f(\rho)\right]
\]
\[
+\left(f'(\rho), -i[H, \rho] -\frac12 \sum_{j=0}^{m-1}\{ L_j^* L_j,\rho\}
+\sum_{j=0}^{k-1} L_j \rho  L_j^* +(m-k) \rho \right)
\]
\begin{equation}
\label{eqmixgenerun}
+\frac12 \sum_{j=0}^{k-1}  \left(\rho  L_j^* +  L_j\rho- {\tr} (\rho L_j^* +L_j \rho)\rho,
f''(\rho)  (\rho  L_j^* +  L_j\rho- {\tr} (\rho L_j^* +L_j \rho) \rho)\right).
\end{equation}

By Ito's formula, the corresponding SDEs yielding an alternative description
of the Markov process generated by \eqref{eqmixgenerun} have the following form
\[
d\rho(t)=\sum_{j=k}^{m-1} (L_j\rho L_j^*-\rho) \, dN_j(t)
\]
\[
+\left( -i[H, \rho] -\frac12 \sum_{j=0}^{m-1}\{ L_j^* L_j,\rho\}
+\sum_{j=0}^{k-1} L_j \rho  L_j^* +(m-k) \rho \right)\, dt
\]
\begin{equation}
\label{Lindstochmix}
+\sum_{j=0}^{k-1} [L_j\rho+\rho L_j^*-\rho\, {\tr} \, (L_j\rho(t)+\rho L_j^*) ] dB_j(t),
\end{equation}
where $N_j(t)$, $j=k, \cdots, m-1$, are independent standard Poisson processes of unit intensity
and $B_j(t)$, $j=0,\cdots, k-1$, are independent (and independent from all $N_j$) standard Brownian motions.

A rigorous construction of the processes generated by
\eqref{eqmixgener}, \eqref{eqjumpgener} , \eqref{eqdifgener}
will be given in Section \ref{secwellpos},
 and their derivation from basic quantum mechanics postulates in Section \ref{secder}.

\subsection{Interaction picture and mild equations}
\label{secinter}

As was mentioned, in this survey paper we decided to reduce the attention mostly
 to the case of bounded coupling operator $L$,
which will be assumed throughout if not stated otherwise.

For the case of bounded $L$ and unbounded $H$ there exists a standard trick,
the interaction representation,
that allows one to always reduce the analysis to bounded (even vanishing) Hamiltonians, though by moving
to time dependent $L$.  Namely, in terms of the new states $\nu(t)=e^{iHt}\ga(t) e^{-iHt}$,
equation \eqref{Lindstoch} rewrites equivalently as the equation
\begin{equation}
\label{Lindstochin}
d\nu(t)=\LC_{L^{Ht}} \nu(t) \, dt +(L^{Ht}\nu(t)+\nu(t) (L^{Ht})^*) dY(t),
\end{equation}
where $L^{Ht}=e^{iHt} L e^{-iHt}$. Similarly, equation \eqref{Lindstochnorm1} rewrites in terms of
$\mu(t)=e^{iHt}\rho(t) e^{-iHt}$ as the equation
\begin{equation}
\label{Lindstochnormin}
d\mu(t)=\LC_{L^{Ht}} \mu (t)\, dt
+[L^{Ht}\mu(t)+\mu(t) (L^{Ht})^*-\mu(t)\, {\tr} \, (L^{Ht}\mu(t)+\mu(t) (L^{Ht})^*) ] dB(t).
\end{equation}

 This rewriting is of course obtained by the trivial application
of Ito's formula that is valid for bounded $H$ and $L$. But
equations  \eqref{Lindstochin} and \eqref{Lindstochnormin} can be looked at as particular cases
of  \eqref{Lindstoch} and \eqref{Lindstochnorm1} with bounded (vanishing ) $H$,
though with time-dependent coefficients even when the initial equation included unbounded $H$.

In future, we shall speak about equation \eqref{Lindstoch} or \eqref{Lindstochnorm1} generally,
though having in mind this exact equation only for bounded $H$ and equation \eqref{Lindstochin}
or \eqref{Lindstochnormin} in general case. Though after this reduction the coupling operator becomes
time-dependent, we shall not usually pay much attention to this detail, because, as mentioned
above in Remark \ref{remtimedep}, all our results modify automatically to treat such time dependence.

Interaction representation is convenient for abstract analysis. For concrete equations
of classical quantum mechanics, it is often handy to use another trick to get rid of unbounded
Hamiltonians, namely, to use the so-called mild (integral) form of stochastic equations:
\[
\ga(t)=e^{-iHt}\ga_0e^{iHt}
+\int_0^t e^{-iH(t-s)} \LC_L\ga (s)  e^{iH(t-s)} \, ds
\]
\begin{equation}
\label{Lindstochmild}
 +\int_0^t e^{-iH(t-s)} (L\ga(s)+\ga(s) L^*) e^{iH(t-s)}\, dY(s)
\end{equation}
for equation \eqref{Lindstoch} and
\[
\rho(t)=e^{-iHt}\rho_0 e^{iHt}
+\int_0^t e^{-iH(t-s)} \LC_L\rho (s)  e^{iH(t-s)} \, ds
\]
\begin{equation}
\label{Lindstochnormmild}
+\int_0^t e^{-iH(t-s)} \bigl[
L\rho(s)+\rho(s) L^*-\rho(s)\, {\tr} \, (L\rho(s)+\rho(s) L^*) \bigr] e^{iH(t-s)} dB(s)
\end{equation}
for equation \eqref{Lindstochnorm1}.

The equivalence of \eqref{Lindstochin} (resp. \eqref{Lindstochnormin}) and \eqref{Lindstochmild}
(resp. \eqref{Lindstochnormmild}) is straightforward again by formal application of Ito's formula.

For the corresponding pure states the mild versions
 of liner and nonlinear quantum filtering equations have the form
\begin{equation}
\label{eqqufiBlinsm}
\chi(t) =e^{-iH t}\chi_0+\int_0^t e^{-iH(t-s)} [-\frac12 L^*L \chi(s) \,dt+L \chi(s) \, dY(s)],
\end{equation}
and, respectively,
\[
\phi(t)=e^{-iH t}\phi_0+\int_0^t e^{-iH (t-s)} \bigl[ i\langle L_S \rangle_{\phi(s)} L_A \phi(s) \, ds
-\frac12 (L-\langle L_S \rangle_{\phi(s)})^*(L-\langle L_S \rangle_{\phi(s)})\phi(s) \, ds
\]
\begin{equation}
\label{eqqufiBnonlinsm}
+ (L-\langle L_S \rangle_{\phi(s)})\phi(s) \, dB(s)\bigr].
\end{equation}

The same tricks, interaction picture and mild forms,
allows one to reduce general filtering equations for counting observation
to the case of vanishing $H$. For instance, equation \eqref{eqBeleqcount} in terms of the states
$\mu(t)=e^{iHt}\rho(t) e^{-iHt}$ rewrites as the equation in interaction picture
 \begin{equation}
\label{eqBeleqcountin}
d\rho=(-\frac12 \{(L^{Ht})^*L^{Ht},\rho\}+ {\tr} (L^{Ht}\rho (L^{Ht})^*) \rho ) \, dt
+\left(\frac{L^{Ht}\rho (L^{Ht})^*}{{\tr} (L^{Ht}\rho (L^{Ht})^*)}-\rho\right) dN(t),
\end{equation}
with the same $L^{Ht}=e^{iHt} L e^{-iHt}$ as above.

\section{Well-posedness}
\label{secwellpos}

\subsection{Preliminaries: equations for pure states}

As a preliminary step in our analysis of filtering equations, we shall briefly look at
filtering equations for pure states reducing attention to the diffusive case only.

If $H$ is also bounded,
equations \eqref{eqqufiBlins} are clearly well-posed by the standard Ito theory
in Hilbert spaces, as they are
linear equations with bounded coefficients.

If $H$ is unbounded, we can use the same trick as for the mixed state equation,
namely, the interaction representation, to get rid of unbounded $H$.
It is seen directly that in terms of vectors
$\xi(t)=e^{iHt} \chi(t)$ and $\psi(t)=e^{iHt} \phi(t)$,
equations \eqref{eqqufiBlins} and \eqref{eqqufiBnonlins},
are equivalent (for bounded $H$) to the SDEs
\begin{equation}
\label{eqqufiBlinsin}
d\xi(t) = -\frac12 (L^{Ht})^*L^{Ht} \xi(t) \,dt+L^{Ht} \xi(t) dY(t),
\end{equation}
and, respectively,
\[
d\psi(t)= i\langle L^{Ht}_S \rangle_{\psi(t)} L^{Ht}_A \psi(t) \, dt
-\frac12 (L^{Ht}-\langle L^{Ht}_S \rangle_{\psi(t)})^*(L^{Ht}-\langle L^{Ht}_S \rangle_{\psi(t)})]\psi(t) \, dt
\]
\begin{equation}
\label{eqqufiBnonlinsin}
+ (L^{Ht}-\langle L^{Ht}_S \rangle_{\psi(t)})\psi(t) \, dB(t),
\end{equation}
with $L^{Ht}=e^{iHt} L e^{-iHt}$.

These SDEs in the interaction form have the same structure as the initial equations, the only difference being
the time dependence of $L$. As for mixed states,
we will thus speak about the well-posedness of linear equations \eqref{eqqufiBlins} meaning its interaction version
\eqref{eqqufiBlinsin} in case of unbounded $H$.

As mentioned above, solution to nonlinear equation  \eqref{eqqufiBnonlinsn} can be obtained by the normalisation
  from a solution of the linear one, which yields, in particular, the existence of a weak solution to this equation.
  On the other hand, the well-posedness of  \eqref{eqqufiBnonlinsn}  was not an obvious task to achieve
  (it was noted in 2008, see\cite{MoraRebo}, that the uniqueness, even of a weak solution,
was still open for bounded $H,L$ in infinite-dimensional spaces). However, as was noted in \cite{KolQuantLLN},
the  well-posedness even in the strong probabilistic sense follows from the standard well-posedness result
  for stochastic Ito's equation with Lipschitz coefficients, due to
  the following observation.

\begin{lemma}
 \label{lemboundder}
 For any bounded operator $M$ in $\HC$, the mapping
 $\HC \to \HC$ given by formulas
 \[
 \psi \to f(\psi)= \langle M \rangle_{\psi} \psi
\]
is continuously Fr\'echet differentiable outside  $\psi= 0$, with the derivative mapping
\[
Df(\psi)[\phi]=(Df(\psi),\phi)
=[(M\hat \psi, \phi)+(M^*\hat \psi, \phi)] \hat \psi
 +(\hat \psi, M\hat \psi)\left[\phi-2\hat \psi \,Re\, (\hat \psi, \phi)\right],
 \]
 where $\hat \psi=\psi/\|\psi\|$. The linear operator $\phi\mapsto Df(\psi)[\phi]$
has the norm not exceeding $3\|M\|$.
In particular, the mapping $f$ is globally Lipschitz with Lipschitz constant $3\|M\|$.
\end{lemma}

\begin{proof}
We have:
\[
\left.\frac{d}{dt}\right|_{t=0} f(\psi+t\phi)
=\left.\frac{d}{dt}\right|_{t=0} \frac{(\psi+t\phi, M(\psi+t\phi))}{(\psi+t\phi, \psi+t\phi)} (\psi+t\phi)
\]
\[
=\left.\frac{d}{dt}\right|_{t=0}\bigl[((\psi,M\psi)+t(\phi, M\psi) +t(\psi, M\phi))
\left(1-\frac{2t \, Re\, (\phi,\psi)}{(\psi,\psi)}\right)\frac{\psi+t\phi}{(\psi,\psi)}\bigr]
=Df(\psi)[\phi],
\]
implying Gateaux differentiability of required form. Since
\[
\|\phi-2\hat \psi \,Re\, (\hat \psi, \phi)\|=\|\phi\|
 \]
(as one sees easily by choosing coordinate  with $\hat \psi$ a coordinate vector),
it follows that the norm of $Df(\psi)$ is bounded by $3\|M\|$. We complete the proof by noting
the obvious continuity of the mapping  $\psi\mapsto Df(\psi)$ from $\HC$ to $\LC(\HC)$,
since the continuity of a Gateaux derivative is known to imply that this derivative is Fr\'echet,
see e.g. Theorem 45 from \cite{Toolbox}.
\end{proof}

\begin{remark} Lipschitz continuity of the function $f$ was proved in \cite{BarchBook}
 for finite-dimensional $\HC$ with a bit more lengthier
calculations. Simple general proof via differentiability is taken from \cite{Kol25a}.
\end{remark}

%Thus it follows the strong well-posedness of both linear and nonlinear
%filtering equations for pure states in case of bounded $L$ (stressing again that for unbounded
%$H$ we mean here actually the interaction representation equations \eqref{eqqufiBlinsin}
%and \eqref{eqqufiBnonlinsin} or the corresponding mild forms).
The link between the two descriptions (linear and normalised) of quantum filtering
is summarised in the following statement.

\begin{prop}
\label{linnorrmpure}

(i) If $\chi(t)$ satisfies \eqref{eqqufiBlins},
then $\|\chi(t)\|^2$ satisfies the equations
 \begin{equation}
\label{chisquare}
d\|\chi(t)\|^2=2(\chi(t),L_S \chi(t))dY(t),
\end{equation}
 \begin{equation}
\label{chisquare1}
d\frac{1}{\|\chi(t)\|^2}=-\frac{2}{\|\chi(t)\|^2}
 \frac{(\chi(t),L_S \chi(t))}{\|\chi(t)\|^2}
\left[dY(t)- \frac{(\chi(t),L_S \chi(t))}{\|\chi(t)\|^2} dt\right],
\end{equation}
and the normalised states $\phi(t)=\chi(t)/\|\chi(t)\|$ satisfy the
nonlinear equation \eqref{eqqufiBnonlins}, where
\begin{equation}
\label{eqdefinnov1}
dB(t)=dY(t)-2(\phi(t), L_S \phi(t)) \, dt.
\end{equation}

(ii) Let $\phi(t)$ have unit norms for all $t$ and satisfy the
nonlinear equation \eqref{eqqufiBnonlins}. Define $\|\chi(t)\|^{-2}$ as the solution
(with the initial condition equal to $1$) to the equation
 \begin{equation}
\label{chisquare4}
d\frac{1}{\|\chi(t)\|^2}=-\frac{2}{\|\chi(t)\|^2}
(\phi(t),L_S \phi(t)) dB_j(t),
\end{equation}
which is seen to be identical with \eqref{chisquare1}
when $B$ and $Y$ are linked via \eqref{eqdefinnov} and $\phi(t)=\chi(t)/\|\chi(t)\|$.
 Then the vectors
$\chi(t)=\phi(t) \|\chi(t)\|$ satisfy the linear equation \eqref{eqqufiBlins}.

(iii) The square norm $\|\chi(t)\|^2$ (respectively its inverse) of a solution
to \eqref{eqqufiBlins} is a positive martingale under
the probability law where $Y(t)$ (resp. $B(t)$) is a Brownian motion. If $B(t)$ is
a Brownian motion and $Y$ is given by \eqref{eqdefinnov}, then
 \begin{equation}
\label{chisquare5}
\E \|\chi(t)\|^2 \le \exp\{2 t \|L\|^2 \} \|\chi_0\|^2,
\end{equation}
where the expectation is with respect to $B(t)$.
\end{prop}


\begin{proof} All formal manipulations are straightforward applications of Ito's lemma and we omit them.
Since well-posedness of all equations is established, these formal manipulations are fully
justified proving (i) and (ii).
Of course, for unbounded $H$ the calculations are performed in the interaction representation.
To prove (iii) we observe from \eqref{chisquare} that $\|\chi\|^2$ is an exponential local martingale solving
the linear equation $d\|\chi\|^2=\|\chi\|^2 dM(t)$ with the martingale
\[
 M(t)=2\int_0^t \left(\frac{\chi(t)}{\|\chi (t)\|},L_S \frac{\chi(t)}{\|\chi (t)\|}\right)dY_j(t).
 \]
 If $Y(t)$ is a BM, this martingale has bounded quadratic variation. Therefore $\|\chi(t)\|^2$ is a true
 (almost sure positive) martingale according to the standard Novikov criterium.
 The same argument applies to $\|\chi(t)\|^{-2}$. Finally, from \eqref{eqdefinnov}, \eqref{chisquare}, it follows that
 \[
 \frac{d}{dt} \E \|\chi(t)\|^2 \le 2\|L\|^2 \, \|\chi(t)\|^2
 \]
 and  \eqref{chisquare5} follows by Gronwall's lemma.
\end{proof}

Let us note that much
more general linear equations in Hilbert spaces preserving the expectation of norm squared,
with bounded coefficients, are treated in detail in \cite{BarchHol}. Nonlinear equations are
more subtle. The well-posedness in strong probabilistic sense for general bounded $L$ was
first obtained in \cite{KolQuantLLN} and for unbounded $L$ in \cite{Kol25b}. In \cite{BarchHol}
 the existence of solutions in a rather general case was proved.
In \cite{MoraRebo} the well-posedness in the weak probabilistic sense was obtained
for unbounded $H,L$ under standard assumptions inspired by the works on the
conservativity of quantum dynamic semigroups from \cite{ChebFagn} and \cite{ChebQuez}.
For other results on the solutions to some generalised
versions of equations \eqref{eqqufiBnonlins} with unbounded $H$ and $L$
under various nontrivial assumptions we can refer also to  \cite{Holevo91}, \cite{Holevo96}, \cite{Fagnola},
\cite{Mora13}. For other classes  of stochastic Schr\"odinger equation we can refer to \cite{BarbRock16},
\cite{GreckGenerNonlinSchr} and references therein.

\subsection{Counting measurements}

Turning to our main topic, mixed states,
we start with the apparently simpler case of counting observation. The
stochastic calculus of jump processes in Banach spaces is more or less straightforward,
because all stochastic integrals can be understood
in the classical Lebesgue sense.
We shall omit details concerning simpler equation \eqref{eqBeleqcountlin}
and will concentrate on the most important (normalised) version \eqref{eqBeleqcount}
considering it in the (physically natural) Banach space
$\HC^1_s$ of self-adjoint trace-class operators.

As a tool, we will need the following counterpart of Lemma \ref{lemboundder} for density operators.

\begin{lemma}
 \label{lemboundderden}
 (i) For any bounded operator $M$ in $\HC$, the mapping
 $\HC^1 \to \HC^1$ given by the formula
 \[
 \ga \to f(\ga)= \frac{{\tr}\, (\ga M)}{{\tr}\, \ga} \ga
\]
is Fr\'echet differentiable outside the hyperplane ${\tr}\, \ga =0$,  with the bounded derivative mapping
\[
Df(\ga)[\de]=(Df(\ga), \de)
=\hat \ga \, {\tr}\, (M\de) +(\de-\hat \ga \,{\tr}\, \de) \,{\tr}\, (\hat \ga M),
 \]
 where $\hat \ga=\ga/{\tr} \, \ga$. (ii) For positive $\ga$, the linear operator $\de\mapsto Df(\ga)[\de]$
in $\HC^1$ has the norm not exceeding $3\|M\|$.
In particular, the mapping $f$ is globally Lipschitz on $S(\HC)$. (iii) The mapping $f(\ga)$ is infinitely differentiable
outside the hyperplane ${\tr}\, \ga =0$, with the derivative of any order bounded on the set $S(\HC))$.
\end{lemma}

\begin{proof}
(i) We have:
\[
\left.\frac{d}{dt}\right|_{t=0} f(\ga+t\de)
=\left.\frac{d}{dt}\right|_{t=0}
\bigl[\frac{{\tr}\, (\ga M)+t\, {\tr}\, (\de M)}{{\tr}\, \ga +t\, {\tr}\, \de}
 (\ga+t\de)\bigr]
\]
\[
=\left.\frac{d}{dt}\right|_{t=0}\bigl[ ({\tr}\, (\ga M)+t\, {\tr}\, (\de M))\frac{1}{{\tr}\, \ga}
\left(1-t \frac{{\tr}\, \de}{{\tr}\, \ga}\right)(\ga +t \de)\bigr]
=(Df(\ga),\de),
\]
as claimed. The differentiability in the sense of Fr\'echet follows from the continuity,
as in Lemma  \ref{lemboundder} above. (ii) It follows from the well known inequality
$\|\rho M\|_{\HC^1}\le \|M\| \, \|\rho\|_{\HC^1}$ for any operators $\rho, M$. (iii)
It is seen that, when differentiating, we always get in the denominator
the powers of ${\tr}\, \ga$.
\end{proof}

Let us stress again that, if not stated otherwise, we always assume that $L$ is bounded. Moreover,
usually we perform calculations for bounded $H$ having in mind that this can always be achieved
via the interaction representation, see Section \ref{secinter}, for an arbitrary self-adjoint Hamiltonian.

\begin{theorem}
\label{JumpWellSDE}
SDE \eqref{eqBeleqcount}, or more precisely \eqref{eqBeleqcountind},
is globally well-posed in $S(\HC)$, so that for any $\rho_0\in S(\HC)$ there exists
a unique solution $\rho(t)\in S(\HC)$ defined for all $t>0$.
\end{theorem}


\begin{proof}
Let us look first at the deterministic part of evolution \eqref{eqBeleqcountind}:
$\dot \rho= b(\rho)$ with
\[
b(\rho)= -i[H, \rho] -\frac12 \{L^*L,\rho\}+ {\tr} (L^*L\rho) \rho.
\]
Since this evolution is explicitly trace preserving only for matrices of unit trace, it is convenient
to look at a modified version of this equation (deterministic part of \eqref{eqBeleqcounttp}):
\[
\dot \rho=\tilde b(\rho)=-i[H, \rho] -\frac12 \{L^*L,\rho\}+ {\tr} (L^*L\rho) \frac{\rho}{{\tr}\, \rho}.
\]
It is seen that the equations with $b(\rho)$ and $\tilde b(\rho)$ coincide
on operators of unit trace, but the equation with $\tilde b(\rho)$ is explicitly
trace preserving for arbitrary trace-class matrices. What is even more important,
according to Lemma \ref{lemboundderden}, the function $\tilde b(\rho)$ is Lipschitz
continuous in the set of positive trace-class operators (or in any neighborhood of
this set). This implies that the ODE $\dot \rho=\tilde b(\rho)$, and hence also
$\dot \rho=b(\rho)$, are well-posed globally as long as solutions start in $S(\HC)$,
if one can show that they always remain in the closed convex set of positive (non-negative) operators.
One can show that they really do it by checking
the general criteria for solutions of ODEs staying inside a closed subset of a Banach space from papers
 \cite{Martin} and \cite{Lakshm}. However, in our case, this conclusion can be derived directly. Namely, if
 $\rho(t)$, started in $S(\HC)$ touches the boundary at some time $t$, this means that there appears
  a vector $v$ such that $(v,\rho(t) v)=0$. But then $(v, b(\rho(t)) v)=(v,\tilde b(\rho(t))v)=0$ meaning
  that this property remains for all future times $t$. Thus $(v,\rho(t)v)$ can never become negative.

\begin{remark} This argument actually shows more, than just preservation of $S(\HC)$
by the deterministic part of evolution \eqref{eqBeleqcountind}. It shows that if
we start away from the boundary of $S(\HC)$, then we shall never touch it (because
the argument above can be used also in the inverse time). And moreover,
the set (actually the subspace) of $v$ such that $(v, \rho(t) v)=0$
(and thus $\rho(t)v=0$ by positivity of $\rho$)
is invariant under the evolution $\dot \rho=b(\rho)$.
\end{remark}

Once the deterministic part of  evolution \eqref{eqBeleqcountind} is constructed,
the full evolution can be built in a standard way by interlacing: trajectory jumps
at times $\tau_j$ of the process $N_t^K$ and moves according to the deterministic
part between the jumps. The jumps are seen to preserve $S(\HC)$. Theorem is proved.
\end{proof}


\begin{theorem}
\label{JumpWellPDE}
The operator  \eqref{eqjumpgener} generates a
$C$-Feller semigroup $T_t$ in $C(S(\HC))$ with the spaces
 $C^1(S(\HC))$ and $C^2(S(\HC))$ being invariant subspaces of the domain,
 and $T_s$ are bounded in these spaces
 uniformly for $s\in [0,t]$ with any $t>0$.
\end{theorem}


\begin{proof}
We start again with the deterministic part of
operator \eqref{eqjumpgener}:
\[
\AC_{count, det}f(\rho)=(f'(\rho), b(\rho)=-(f'(\rho), i[H, \rho]
+ \frac12 \{L^*L,\rho\}-\rho \, {\tr}\, (L^*L\rho)).
\]
As was shown at the beginning of the proof of Theorem \ref{JumpWellSDE} above the ODE
$\dot \rho=b(\rho)$ is globally well-posed in $S(\HC)$. Moreover, as the coefficients
of this equation are infinitely smooth, its solutions
depend infinitely smooth on the initial condition. Hence, by the standard link between ODEs and
first order PDEs (see e.g. details of this link for Banach space in Section 2.10 \cite{Kolbook19}),
resolving operators of the Cauchy problem for equation $\dot f=\AC_{count,det}f$ define the semigroup
$T_t^{det}$ of contractions in the space $C(S(\HC))$ such that all spaces $C^k(S(\HC))$ are invariant and
$T_t$ act as a bounded semigroup in each of this spaces.  Next, it is straightforward to see that
the jump part $\AC_{jump}=\AC_{count, det}-\AC_{count}$ of the operator $\AC_{count}$ represents a
bounded operator in $C(S(\HC))$ and in each space $C^k(S(\HC))$ (we need only $k=1,2$). For instance,
the latter can be seen by finding the derivatives explicitly. For instance,
\[
D\AC_{jump}f(\rho)[\de]
= {\tr}\, (L^*L\de)\left[f(\frac{L\rho L^*}{{\tr}\, (L^*L\rho)})-f(\rho)\right]
\]
\[
+{\tr}\, (L^*L\rho)\left(Df(\frac{L\rho L^*}{{\tr}\, (L^*L\rho)})
\left[ \frac{L\de L^*}{{\tr}\, (L^*L\rho)}-\frac{L\rho L^*}{{\tr}\, (L^*L\rho)}\,
\frac{{\tr}\, (L^*L\de)}{{\tr}\, (L^*L\rho)} \right]
-Df(\rho)[\de]\right).
\]

Since $\AC_{jump}$ is bounded in all spaces $C^k(S(\HC))$, where $\AC_{count,det}$ generates a bounded semigroup,
it follows that $\AC_{count}$ generates there a bounded semigroup $T_t$ as well, due to the basic perturbation theory for
semigroups (see e.g. \cite{ReedSimon} or Sec 4.6 of \cite{Kolbook19})). It is seen directly that the operators $T_t$
in the space $C(S(\HC))$ act as positivity preserving contractions (of course, it follows also from
the fact that the operator $\AC_{jump}$ is conditionally positive).
Finally, the fact that  \eqref{eqjumpgener} generates the process given
by SDEs  \eqref{eqBeleqcount} follows from Ito's formula.
\end{proof}

For the case of finite-dimensional $\HC$ it follows that the spaces $C^k(S(\HC))$
are dense in $C(S(\HC))$ and we can conclude that $C^1(S(\HC))$ and $C^2(S(\HC))$ represent invariant
cores for \eqref{eqBeleqcount} yielding the proper analytic characterisation of the semigroup $T_t$.
In infinite-dimensional case an annoying detail does not allow for a similar conclusion.
It turns out to be a very hard problem to assess the richness of the class $C^1(S(\HC))$ of smooth functions.
Seemingly, it is not dense in $C(S(\HC))$, see lengthy discussions of similar questions in
\cite{Toolbox} and \cite{Phelps}.

Therefore, to complete the analytic characterization of  $T_t$ for infinite-dimensional $\HC$
we supplement Theorem \ref{JumpWellPDE} with the following claim:

\begin{theorem}
\label{JumpWellPDE1}
An invariant core for the generator of $T_t$ can be chosen as the space
$C^1_{Gat*}(S(\HC))$ consisting of continuous uniformly Gateaux differentiable functions on $S(\HC)$
with the mapping $\ga\mapsto f'(\ga)=Df(\ga)$
from $S(\HC)$ to $(\HC^1_s)^*$ being bounded and continuous
when $S(\HC)$ is equipped with its norm topology
and $(\HC^1_s)^*$ with its weak* topology.
\end{theorem}

\begin{proof}
The fact that $C^1_{Gat*}(S(\HC))$ is an invariant subspace of the domain of \eqref{eqBeleqcount}
is obtained as above. Notice only that the required continuity is sufficient to ensure that
$L_{count}f \in C(S(\HC))$ for any $f\in C^1_{Gat*}(S(\HC))$. The nontrivial fact is that
$C^1_{Gat*}(S(\HC))$ is dense in $C(S(\HC)$. To see that this is true we first note that,
similar to the theory of $l^1$ spaces (Theorem 5.12 of \cite{Phelps}), one can show the existence
of an equivalent Gateaux differentiable norm on $\HC^1_s$. Equivalently one can refer to a more
result, see \cite{Hayek}, that in any separable Banach space there exists a uniformly
Gateaux differentiable equivalent norm.
Next, by Proposition 2.8 of \cite{Phelps}, the derivative of such a norm is norm-weak${}^*$
continuous mapping from $\HC^1_s$ to $(\HC^1_s)^*$. Finally, by the well known theory
(see e.g. \cite{Eells}), the existence of norms of certain smooth class in a separable
Banach space implies the possibility of approximating continuous functions by the
smooth functions of the same class.
\end{proof}

If the coupling operator $L$ is unitary yielding the linear filtering equation
\eqref{eqBeleqcountunit}, a more straightforward and a more detailed analytic
characterization of the core for the generator of $T_t$ is available (yielding
cores of arbitrary smoothness classes)
by two different approaches.
The first one is given by the following.

\begin{prop}
\label{JumpWellPDE2}
If $L$ is unitary, the semigroup  $T_t$ preserves the subspace of bounded weakly* continuous functions
on $S(\HC)$, and their subspaces of smooth functions of first or second order are invariant cores
for the generator of so restricted $T_t$.
\end{prop}

\begin{proof} The preservation of weakly* continuous functions by $T_t$ follows from linearity
of \eqref{eqBeleqcountunit} (which is not at all obvious for the general case). The fact that smooth
functions are dense in the space of weakly* continuous functions follows from Stone-Weierstrasse theorem.
\end{proof}

The second method can be developed via the Hilbert-space lifting of the dynamics of \eqref{eqBeleqcountunit}.
It will be presented below for more general mixed channels.


\subsection{Diffusive case: linear version}
\label{seclineq}

The first thing to decide for dealing with the equations on mixed states is the choice
of an appropriate Banach space of operators, where the corresponding SDEs will be analysed.

For the diffusive case, we shall consider these equations in the Hilbert space $\HC^2_s$ of self-adjoint Hilbert-Schmidt
operators in $\HC$, which is a closed subspace
in the space of all Hilbert-Schmidt operators $\HC^2$ with the scalar product
${\tr} (A^*B)={\tr} (AB)$. Note that $\HC^2_s$ is a real Hilbert space
(though consisting of operators in the complex Hilbert space $\HC$).

%\begin{remark}
Since we are interested in trace-class operators, a more natural
space from physical point of view  would be of course the Banach space $\HC^1_s$ of self-adjoint
trace-class operators in $\HC$ (that we used above for the counting case).
However, the classes of Banach spaces,
for which a satisfactory extension of Ito stochastic calculus was developed,
namely the so-called UMD spaces, spaces of martingale type 2 and spaces with a smooth norm
(see review \cite{VanNeerven}) do not include $\HC^1$, and therefore we work in the larger space
$\HC^2$. In this space the key linear functional of taking trace is unbounded,
and we are led to work with SDEs with singular coefficients. This complication
is the price to pay for working in a convenient Hilbert setting of $\HC^2$.
%\end{remark}



\begin{theorem}
\label{LindStochLin1}
Assume $Y(t)$ is a ($n$-dimensional) standard BM.
 Then the following claims hold:

(i) Equation \eqref{Lindstoch} is well-posed in $\HC^2_s$, that is,
it has a unique global solution for any $\ga_0\in \HC^2_s$.
These solutions have the growth estimates
 \begin{equation}
\label{Lindstochquadex1}
\E [{\tr}\, \ga^2(t)] \le {\tr}\, \ga_0^2 \exp\{ 4 t\|L\|^2\}.
\end{equation}

(ii) The solution $\ga (t)$ to \eqref{Lindstoch}
is positive-definite for all $t$ whenever $\ga_0$ so is.

(iii) If the initial condition $\ga_0$ is of trace-class,
then so is the solution $\ga(t)$, with the trace given by the formula
\begin{equation}
\label{Lindstochmart}
{\tr} \,\ga(t)={\tr} \, \ga_0+\int_0^t {\tr} (L\ga(s)+\ga(s) L^*) dY(s).
\end{equation}
Moreover, ${\tr} \, \ga(t)$ is a
square integrable martingale such that
\begin{equation}
\label{Lindstochmart1}
\E ({\tr} \,\ga(t))^2\le [({\tr} \, \ga_0^+)^2 +({\tr} \, \ga_0^-)^2]\exp\{ 4t \|L\|^2\},
\end{equation}
where $\ga_0^{\pm}$ denote positive and negative parts of $\ga_0$, and
\begin{equation}
\label{Lindstochmart2}
\E ({\tr} \,|\ga(t)|)\le {\tr} \, |\ga_0|.
\end{equation}

\end{theorem}

%\begin{remark}
%The proof of positivity of the solutions was usually considered as a difficult task.
%Three different approaches for such proof in finite-dimensional case can be found e.g. in
%\cite{BarchBook} (using a link with pure state equations), \cite{Pellegrini10} (using discrete
%Markov  chain approximation), \cite{KolQuantFrac} (turning to Stratonovich SDE or using the theory
%of attainable boundary points). We give here a simple general proof from \cite{Kol25a}.
%\end{remark}

\begin{proof}
(i) The existence of a unique solution is straightforward, as the coefficients at $dt$ and $dY$ are
bounded linear operator of $\ga$ in $\HC^2_s$.

We shall work with equation \eqref{Lindstoch}, extension to  \eqref{Lindstochin} being automatic.

We derive from
\eqref{Lindstoch} by Ito's formula that
\begin{equation}
\label{Lindstochquadtr}
d \, {\tr}\, \ga^2={\tr} \,(2\ga L\ga L^*+ L \ga L \ga + \ga L^* \ga L^*) \, dt
+2\, {\tr} \, [(L+L^*)\ga^2] \, dY(t),
\end{equation}
and thus
\[
\E \, {\tr}\, \ga^2(t)
={\tr}\, \ga_0^2 +\E \, {\tr} \,
\int_0^t (\ga(s) L^*\ga(s) L+\ga(s) L\ga(s) L^*+ L \ga(s) L \ga(s) + \ga(s) L^* \ga(s) L^*) ds
\]
so that, by \eqref{eqmyineqtr1},
\[
\E \, {\tr}\, \ga^2(t) \le {\tr}\, \ga_0^2+ 4 \|L\|^2 \int_0^t  \E \, {\tr} \,(\ga^2(s)) \, ds.
\]
and \eqref{Lindstochquadex1} follows by By Gronwall's lemma.

 (ii) Since $\ga_0$ is a positive operator from $\HC^2_s$, it follows that
  there exists a orthonormal basis $\{e_k\}$ in $\HC$ such that $\ga_0$
 can be presented as a convergent (in $\HC^2_s$) series
 \[
 \ga_0=\sum_{k=1}^{\infty} p_k e_k\otimes \bar e_k
 \]
 with non-increasing non-negative sequence $\{p_k\}$ from $l^2$. Hence $\ga_0=\lim \ga_{0n}$
 with finite-dimensional operators
 \[
 \ga_{0n}=\sum_{k=1}^n p_k e_k\otimes \bar e_k.
 \]
By linearity (and uniqueness of solutions), the solution $\ga_n(t)$ with the initial condition $\ga_{0n}$ is the finite
convex combination of the solutions with the initial conditions $e_k\otimes \bar e_k$, the latter
being given by $e_k(t)\otimes \bar e_k(t)$ with $e_k(t)$ solving the linear filtering equation for
pure states \eqref{eqqufiBlins}, and thus being positive definite. Therefore, $\ga_n(t)$ are positive-definite.
On the other hand, by \eqref{Lindstochquadex1}, $\ga_n(t)-\ga(t)$ tend to zero, as $n\to \infty$,
because the solutions with initial condition $\ga_0-\ga_{0t}$ tend to zero. Hence all $\ga (t)$ are
also positive-definite.

(iii) First assume that $\ga_0$ is positive. We then use the approximations
$\ga_n(t)$ as defined in (ii) above.
Since all $e_k(t)\otimes \bar e_k(t)$ are positive-definite operators,
the sequence $\ga_n(t)$ is monotonically increasing in $n$.

Since
\begin{equation}
\label{tracenormsq}
{\tr} \, e_k(t)\otimes \bar e_k(t)=\|e_k(t)\|^2,
\end{equation}
it follows that, for $n>m$,
\[
{\tr}\, |\ga_n(t)-\ga_m(t)|={\tr}\, (\ga_n(t)-\ga_m(t))
=\sum_{k=m+1}^n p_k \|e_k(t)\|^2,
\]
and thus the sequence $\ga_n(t)$ converges not only in $\HC^2$, but also in
the space of trace-class operators $\HC^1$. Hence, $\ga(t)$ is of trace class and
${\tr }\, \ga(t)=\lim {\tr}\, \ga_n(t)$.

From \eqref{chisquare} and \eqref{tracenormsq} it follows that
\[
{\tr} \,\ga_n(t)={\tr} \, \ga_{0n}+\int_0^t {\tr} (L\ga_n(s)+\ga_n(s) L^*) dY(s).
\]
Passing to the limit in this equation
we obtain \eqref{Lindstochmart}.

By \eqref{Lindstochmart}, Ito's formula and the estimate
$|{\tr} (\ga L)|\le {\tr}\, \ga \,  \|L\|$, it follows that
 \[
 \E ({\tr}\, \ga(t))^2
 \le  ({\tr} \, \ga(0))^2+4 \|L\|^2 \int_0^t \E ({\tr} \, \ga(s))^2\,  ds
\]
implying \eqref{Lindstochmart1} by Gronwall's lemma.

Finally, if $\ga_0$ is not positive, we decompose $\ga_0$
as the difference of its positive and negative parts:
$\ga_0=\ga_0^+-\ga_0^-$. Applying  \eqref{Lindstochmart}
 to the corresponding solutions $\ga^{\pm}(t)$ , we get  \eqref{Lindstochmart} for $\ga(t)$
by linearity. Similarly estimates \eqref{Lindstochmart1} are obtained. It remains estimate \eqref{Lindstochmart2}.
Firstly it clearly holds with the sign of equality for positive $\ga_0$. For general Hermitian $\ga_0$,
we can write
\[
{\tr} \, |\ga(t)| \le {\tr} \, \ga^+(t) +{\tr} \, \ga^-(t).
\]
The sign of inequality is due to the fact that, though the solutions $\ga^{\pm}(t)$ are positive operators,
 they are not necessarily positive and negative parts of the operator $\ga(t)$).
Applying  \eqref{Lindstochmart2} to $\ga^{\pm}(t)$ we get \eqref{Lindstochmart2} to $\ga(t)$.
 \end{proof}

Important part of well-posedness of an equation is the continuous dependence of its solution on
 initial conditions and parameters. By linearity, continuous dependence on initial condition
 follows from \eqref{Lindstochquadex1}. Next result establishes the continuous dependence on
 a bounded part of the Hamiltonian.

\begin{theorem}
\label{LindStochLin2}
Under the assumptions of Theorem \ref{LindStochLin1},
consider two equations of type \eqref{Lindstoch}:
\begin{equation}
\label{Lindstochtwo}
d\ga(t)=-i[H,\ga] \, dt -i[H_j,\ga] \, dt+\LC_L \ga(t) \, dt +(L\ga(t)+\ga(t) L^*) dY(t),
\end{equation}
$j=1,2$, where $H_1$ and $H_2$ are two bounded self-adjoint operators in $\HC$.
Then for their solutions $\ga_j(t)$, $j=1,2$, with one and the same positive initial
condition $\ga_0$ of trace-class
we have the estimates for the deviations in the norms of $\HC^1$ and $\HC^2$:
\begin{equation}
\label{eqLindstochLin21}
\E \, {\tr}\, |\ga_1(t)-\ga_2(t)|\le 2 t \|H_2-H_1\| \, {\tr} \, \ga_0,
\end{equation}
\begin{equation}
\label{eqLindstochLin22}
\sqrt{\E \, {\tr}\, (\ga_1(t)-\ga_2(t))^2}
\le 2 t \|H_2-H_1\|\sqrt{ {\tr} \, \ga_0^2} \exp\{2t \|L\|^2\}.
\end{equation}
\end{theorem}

\begin{proof}
By subtracting the two equations we get the following:
\[
d(\ga_1(t)-\ga_2(t))=-i[H_1,\ga_1(t)-\ga_2(t)] \, dt
+\LC_L (\ga_1(t)-\ga_2(t)) \, dt +(L(\ga_1(t)-\ga_2(t))
\]
\begin{equation}
\label{Lindstochtwo1}
+(\ga_1(t)-\ga_2(t)) L^*) dY(t)
+i[H_2-H_1, \ga_2(t)] \, dt.
\end{equation}
Denoting by $\Phi_t$ the operator giving solution to the Cauchy problems for equation \eqref{Lindstochtwo}
with $j=1$, we can express solution to \eqref{Lindstochtwo1} with the vanishing
initial condition in the following standard Du Hamel form:
\begin{equation}
\label{Lindstochtwo2}
\ga_1(t)-\ga_2(t)=i\int_0^t \Phi_{t-s} [H_2-H_1, \ga_2(s)] \, ds.
\end{equation}
Hence
\[
\E \, {\tr}\, |\ga_1(t)-\ga_2(t)|\le \int_0^t \E \, {\tr} \, |\Phi_{t-s} [H_2-H_1, \ga_2(s)]| \, ds.
\]
By the chain rule,  one can insert the conditional expectation with
respect to $\FC_s$ inside the expectation on the r.h.s. of the inequality
and then apply \eqref{Lindstochmart2} leading to the estimate
\[
\E \, {\tr}\, |\ga_1(t)-\ga_2(t)|\le \int_0^t \E \, {\tr} \, | [H_2-H_1, \ga_2(s)]| \, ds
\]
\[
\le 2\|H_2-H_1\| \int_0^t \E \, {\tr} \, \ga_2(s) \, ds
\le 2 t \|H_2-H_1\|  \, {\tr} \, \ga_0,
\]
implying \eqref{eqLindstochLin21}.

Similarly, using \eqref{Lindstochtwo2}, insertion of conditional expectation
and estimates  \eqref{Lindstochquadex1} and \eqref{eqmyineqtr2}, we write
\[
\sqrt{\E \, {\tr}\, (\ga_1(t)-\ga_2(t))^2}
\le \int_0^t \sqrt{ \E \, {\tr} \, (\Phi_{t-s} [H_2-H_1, \ga_2(s)])^2} \, ds
\]
\[
\le \int_0^t \sqrt{\E \, {\tr} \, ([H_2-H_1, \ga_2(s)])^2} \exp\{2(t-s)\|L\|^2\} \, ds
\]
\[
\le 2\|H_2-H_1\| \int_0^t \sqrt{\E \, {\tr} \, (\ga_2(s))^2} \exp\{2(t-s)\|L\|^2\} \, ds
\]
\[
\le 2 \|H_2-H_1\| \sqrt{{\tr} \, \ga_0^2} \int_0^t \exp\{2t\|L\|^2\} \, ds
\]
implying \eqref{eqLindstochLin22}.
\end{proof}

\begin{remark} An alternative proof of \eqref{eqLindstochLin22} can be given by writing
down the equation for ${\tr} \, (\ga_1(t)-\ga_2(t))^2$ and then applying Gronwall's lemma.
\end{remark}

\subsection{Semigroups of linear filtering equations in diffusive case}

Let us look at the analytic properties of the semigroup $T_tf(\ga)=\E f(\Ga_t \ga)$,
where $\Ga_t$ are random linear operators yielding the solution
to the Cauchy problem of SDE \eqref{Lindstoch} (given by Theorem \ref{LindStochLin1}).
Since $\Ga_t(\ga)$ are continuous, the semigroup $T_t$ is a contraction semigroup in $C(\HC^2_s)$.
It also preserves positivity, and therefore, it is a $C$-Feller semigroup.

\begin{theorem}
\label{Linsdesem}

(i) $T_t$ preserve the spaces $C^k(\HC^2_s)$ for all $k$ and act there as a bounded quasi-contraction semigroup,
that is, their norm is bounded by $e^{ct}$ with a constant $c$.

(ii) If $f\in C_{luc}(\HC^2_s)$, then $T_tf(\ga) \to f(\ga)$, as $t\to 0$,
uniformly on $\ga$ from any bounded set. The space $C_{luc}(\HC^2_s)$ is invariant under $T_t$.

(iii) The space $C_{\infty}(\HC^2_s)$ is invariant under $T_t$.
If $f\in C_{luc}(\HC^2_s)\cap C_{\infty}(\HC^2_s)$,
then $T_tf(\ga) \to f(\ga)$, as $t\to 0$, uniformly in all $\ga$.

(iv) The semigroup $T_t$ is strongly continuous in the closed subspace
$C_{luc, \infty}(\HC^2_s)= C_{luc}(\HC^2_s)\cap C_{\infty}(\HC^2_s)$
of $C(\HC^2_s)$.
\end{theorem}

\begin{proof}
(i) Assuming $f\in C^1(\HC^2_s)$, we have
\[
T_tf(\ga+\rho)=\E f(\Ga_t(\ga+\rho))=\E f(\Ga_t \ga+\Ga_t\rho)
\]
\[
=T_tf(\ga)+ \E (f'(\Ga_t\ga), \Ga_t \rho)+\E \, o(\|\Ga_t\rho)\|
=T_t f(\ga)+\E (\Ga_t^* f'(\Ga_t \ga), \rho)+o(\|\rho\|).
\]
Consequently, $(T_tf)'=\E \Ga_t^* f'(\Ga_t\ga)$ and,
by \eqref{Lindstochquadex1},
\[
\|T_tf\|_{C^1(\HC^2_s)}\le e^{2\|L\|t} \|f\|_{C^1(\HC^2_s)}
\]
implying the claim for $k=1$.
Similarly
\begin{equation}
\label{eqLinsdesem}
(T_tf)''=\E \Ga_t^* f''(\Ga_t\ga) \Ga_t
\end{equation}
and other cases follow analogously.

(ii) Let  $f\in C_{luc}(\HC^2_s)$, and let $B_R$ denote the ball of radius $R$ in $\HC^2_s$.
By \eqref{Lindstochquadex1} and \eqref{Lindstoch}, for any bounded $M\subset \HC^2_s$ and any $T$,
there exists a constant $C(T)$ such that $\E \|\Ga_t \ga \|_{\HC^2_s}\le C(T)$ and
\[
\E \|\Ga_t\ga-\ga\|_{\HC^2_s}\le t C(T),
\]
\[
\E \|\Ga_t\ga_1-\Ga_t \ga_2\|_{\HC^2_s}\le \|\ga_1-\ga_2\| C(T),
\]
for all $t\le T$ and all solutions with $\ga,\ga_1,\ga_2 \in M$.
Hence, by the Markov inequality, for any $R$,
 $\P (\Ga_t\ga \notin B_R)<C(T)/R$, and
\[
\P (\|\Ga_t\ga-\ga\|_{\HC^2_s}>\sqrt t) \le \sqrt t C(T),
\]
\[
\P (\|\Ga_t\ga_1-\Ga_t\ga_2\|_{\HC^2_s}>\sqrt{\|\ga_1-\ga_2\|}) \le \sqrt {\|\ga_1-\ga_2\|}  C(T).
\]

Next, by local continuity of $f$, there exists $\de$ such that
$\|\ga_1-\ga_2\|_{\HC^2_s}\le\de $ and $\ga_1,\ga_2\in B_R$ imply
that $|f(\ga_1)-f(\ga_2)|<1/R$. Consequently, for $\sqrt t <\de$,
\[
\sup_{\ga\in M} |\E f(\Ga_t \ga)-f(\ga)|
\le 2\|f\| (\frac{C(T)}{R}+\sqrt t C(T)) +\frac{1}{R},
\]
which can be made arbitrary small by choosing large $R$ and small $t$. This proves the first claim.

On the other hand, for $\sqrt{\|\ga_1-\ga_2\|}<\de$,
\[
|\E f(\Ga_t \ga_1)-\E f(\Ga_t \ga_2)|\le  \E |f(\Ga_t \ga_1)-f(\Ga_t \ga_2)|
\le 2\|f\| (2 \frac{C(T)}{R}+\sqrt {\|\ga_1-\ga_2\|}  C(T)) +\frac{1}{R},
\]
which can be made arbitrary small by choosing large $R$ and small $\|\ga_1-\ga_2\|$. This proves the second claim.

(iii) The invariance of $ C_{\infty}(\HC^2_s)$ follows from \eqref{Lindstochquadex1}.
If $f\in C_{\infty}(\HC^2_s)\cap C_{luc}(\HC^2_s)$, then, using \eqref{Lindstochquadex1},
we can make all values $T_tf(\ga)$ small for small enough $t$ and $\|\ga\|_{\HC^2_s}>R$
with sufficiently large $R$, and then apply (ii) for $M=B_R$.

(iv) It is clear that $C_{luc}(\HC^2_s)\cap C_{\infty}(\HC^2_s)$
is closed in $C(\HC^2_s)$. The strong continuity is a consequence of (iii).
\end{proof}

Knowing the conservation of smoothness by $T_t$, we can justify the application of Ito's formula,
and conclude rigorously
that for any $f\in C^2(\HC^2_s)$,
the function $T_tf(\ga)=f(t,\ga)$ solves the Cauchy problem
\eqref{eqCauchylinfilt} with the diffusion operator $\AC$ given by
\eqref{eqgenlinfilt}.

However, bounded functions are neither sufficient nor natural in our quantum mechanics setting,
where quadratic or linear functions are mostly met. To deal with them notice firstly that
linear functions of the type $\hat \phi: \ga \mapsto (\phi, \ga)={\tr}\, (\phi \ga)$ with a $\phi \in \HC^2_s$
are invariant under $T_t$. Namely,
\[
T_t \hat \phi (\ga)=\E (\phi, \Ga_t(\ga))=(\E \Ga_t^* \phi, \ga),
\]
where $\Ga_t^*$ is the resolving operator for the dual equation to \eqref{Lindstoch}:
\begin{equation}
\label{Lindstochdu}
dm(t)=i[H,m(t)] \, dt +\LC_{L^*} m(t) \, dt +(L^* m(t)+ m(t) L) dY(t).
\end{equation}

A natural space for $T_t$ is a space $C_{quad}(\HC^2_s)$ of continuous functions $f$
of at most quadratic growth, that is, when the norm
\[
\|f\|_{quad}=\sup_{\ga \in \HC^2_s} \frac{|f(\ga)|}{1+{\tr} \, \ga^2}
\]
is bounded. A convenient subspace of this space is the space $C^2_{quad}(\HC^2_s)$
of twice continuously differentiable functions with bounded second derivatives.
Equipped with the norm
\[
\|f\|_{C^2_{quad}(\HC^2_s)}=\|f\|_{quad}+\sup_{\ga} \frac{\|f'(\ga)\|}{\sqrt{1+{\tr} \, \ga^2}}
+\sup_{\ga} \|f''(\ga)\|,
\]
this space is Banach and such that the generator $\AC$ is bounded as an operator
from $C^2_{quad}(\HC^2_s)$ to $C_{quad}(\HC^2_s)$.

\begin{theorem}
\label{Linsdesem1}

(i) $T_t$ preserve the spaces $C_{quad}(\HC^2_s)$ and $C^2_{quad}(\HC^2_s)$
and act there as a bounded quasi-contraction semigroup.

(ii) The semigroup $T_t$ is strongly continuous in the closed subspace $C_{quad,uc,\infty}(\HC^2_s)$
of $C_{quad}(\HC^2_s)$ consisting of functions $f$ such that the function $f(\ga)/(1+{\tr}\, \ga^2)$
belongs to  $C_{luc,\infty}(\HC^2_s)$.

(iii) If $f\in C^2_{quad}(\HC^2_s)$, then$T_tf$ solves the Cauchy problem \eqref{eqCauchylinfilt}.

(iv)) The subspace of $C_{quad}^2(\HC^2_s)$ consisting of functions with the second derivative from
$C_{luc,\infty}(\HC^2_s)$ is an invariant core for the generator of the semigroup $T_t$ in $C_{quad,uc,\infty}(\HC^2_s)$.
 \end{theorem}

\begin{proof}
(i) The invariance of $C_{quad}^2(\HC^2_s)$ follows from
\eqref{eqLinsdesem}. The preservation of the quadratic growth
 follows from \eqref{Lindstoch}. Part (ii) is a consequence of (i) and Theorem \ref{Linsdesem}.
 Part (iii) is a consequence of Ito's lemma. (iv) Taking into account all previous statements,
it remains to note (see e.g. \cite{Eells}) that smooth functions are dense in the space of continuous functions
on any open bounded subset of a Hilbert space.
\end{proof}

\begin{remark}
The statements of Theorem \ref{Linsdesem1} remain true, if instead of
the space $C_{quad}(\HC^2_s)$ one works with its subspace $C_{quad, weak}^2(\HC^2_s)$ of weakly continuous functions.
This space is preserved by $T_t$, because the solutions of  SDE \eqref{Lindstoch} are given by continuous linear operators.
\end{remark}

%In this case strong continuity of $T_t$ holds in the subspace $C_{quad,weak,\infty}(\HC^2_s)$ of functions satisfying
% \eqref{eqLinsdesem1} (the requirement of uniform continuity is not necessary anymore) and the corresponding subspace of
% $C_{quad}^2(\HC^2_s)$ becomes an invariant core.  The point to note here is that the cylindrical functions of type
%% $g((m_1,\ga), \cdots, (m_k,\ga))$ with continuous bounded $g$ and $m_j\in \HC^2_s$, and therefore also smooth functions
%  of this kind,  are dense in the space of weakly continuous functions on the unit ball in $\HC^2_s$ by
%   the Stone-Weierstrass theorem and due to the weak compactness of this unit ball.
%\end{remark}

\subsection{Diffusive case: normalized version}
\label{secnormeq}

 Recall that a strong
solution of an SDE like \eqref{Lindstochnorm1} is a process $\rho(t)$ that can
 be expressed as a measurable function of a given Brownian motion $B(t)$. By a
weak solution one means a pair of processes $(\rho(t),B(t))$ (where $B(t)$ is a Brownian
motion) defined on a certain stochastic basis, adapted to its filtration and satisfying
\eqref{Lindstochnorm1}. One says that weak solution is  unique in law
if for any two solutions $(\rho^1,B^1)$ and $(\rho^2,B^2)$ (possibly defined on
different probability spaces) the processes $\rho^1$ and $\rho^2$ have the same distribution.


As was mentioned, solutions of nonlinear equation can be built
formally from the linear one. Namely, from \eqref{Lindstochmart}
we can derive by Ito's formula that
\[
d\frac{1}{{\tr} \, \ga(t)}=-\frac{1}{({\tr} \, \ga(t))^2} \, {\tr} \, (L\ga(t) +\ga(t) L^*) dY_t
+\frac{1}{({\tr} \, \ga(t))^3} [{\tr} \, (L\ga(t)+\ga(t) L^*)]^2 \, dt.
\]
Hence by Ito's product rule we check that the normalised density operator
$\rho(t)=\ga(t)/{\tr} \, \ga(t)$ satisfies the equation
 \[
d\rho=-i[H, \rho(t)] \, dt +\LC_L \rho(t)\, dt
\]
 \begin{equation}
\label{Lindstochnorm}
+(L\rho(t)+\rho(t) L^*-\rho(t)\, {\tr} \, (L\rho(t)+\rho(t) L^*) )
[dY_t-{\tr} \, (L\rho(t)+\rho(t) L^*) dt].
\end{equation}

Therefore, in terms of the {\it innovation process}
 \begin{equation}
\label{outputinnovation}
B(t)=Y(t)-\int_0^t {\tr} \, (L\rho(s)+\rho(s) L^*) \, ds
\end{equation}
 the equation for the inverse trace rewrites as
  \begin{equation}
\label{eqtrinnov}
d\frac{1}{{\tr} \, \ga(t)}
=-\frac{1}{{\tr} \, \ga(t)} \, {\tr} \, (L\rho(t) +\rho(t) L^*) dB(t)
\end{equation}
and the equation for the normalized density operator \eqref{Lindstochnorm} rewrites in the
standard form \eqref{Lindstochnorm1} of the nonlinear filtering equation.
 Notice that, for positive solutions, ${\tr} \, \ga(t)$ is a positive martingale
 that specifies the density between the measures $\P$ and $\Q$ that make $Y(t)$
 and $B(t)$ Brownian motions, respectively:
\[
\E_{\Q}\xi=\E_{\P}(\xi \, {\tr}\, \ga(t))
\]
for $\FC_t$-adapted $\xi$.
Therefore, Theorem \ref{LindStochLin1} implies the existence of
weak solutions to nonlinear equation \eqref{Lindstochnorm1}.

Moreover, let a continuous pair of processes $(\rho(t), B(t))$ be a
weak solution to \eqref{Lindstochnorm1} with a positive initial condition
$\rho_0$ and with ${\tr}\, \rho(t)=1$ for all $t$. Then,
if the process ${\tr }\, \ga(t)$ is defined as the solution of equation \eqref{eqtrinnov}
(with any positive initial condition), the process
$\ga(t)$ is defined as $\ga(t)= \rho(t) {\tr }\, \ga(t)$ and the process
$Y(t)$ is defined via \eqref{outputinnovation}, then these processes satisfy
the linear equation \eqref{Lindstoch}. Therefore, any weak solution of the
nonlinear problem can be obtained by normalization from a solution to the linear one.
Hence the uniqueness for the latter implies the uniqueness (in distribution) for the former.

Summarizing we can conclude the following well-posedness
result for weak solutions of nonlinear equation \eqref{Lindstochnorm1}.

\begin{theorem}
\label{LindStochnonLin}
Let $\rho_0$ be a positive operator of unit trace.
Then there exists a unique in law weak solution of equation \eqref{Lindstochnorm1} in $\HC^2_s$
with $B(t)$ a Brownian motion, the initial data $\rho_0$ and such that all $\rho(t)$
are positive-definite operators of unit trace.
\end{theorem}

Let us turn to the strong solutions of \eqref{Lindstochnorm1}. The difficulty with its direct analysis
is due to the fact that though, according to Lemma \ref{lemboundderden}, it can be written
as an equation with Lipschitz coefficients in $\HC^1_s$, the stochastic integrals are defined in $\HC^2_s$,
where the coefficients are singular.
To overcome this difficulty we employ an additional idea:
it turns out that the equation for mixed states can be rewritten as a filtering equation for pure
 states in some enhanced Hilbert space.

 \begin{theorem}
\label{LindStochnonLin1}
Let $\rho_0$ be a positive operator of unit trace,
and $B(t)$ a BM. Then there exists a unique strong solution of equation
\eqref{Lindstochnorm1} in $\HC^2_s$, with the initial data $\rho_0$
and such that all $\rho(t)$ are positive trace class operators of unit trace.
\end{theorem}

\begin{proof}
From the proof of Theorem \ref{LindStochnonLin} we know that $\rho(t)$ solves \eqref{Lindstochnorm1}
 if and only if $\ga(t)={\tr}\, \ga(t)\, \rho(t)$ solves the equation
  \begin{equation}
\label{Lindstochnormlinmix}
d\ga(t)=-i[H, \ga(t)] dt +\LC_L \ga(t) dt +(L\ga(t)+\ga(t) L^*)
[dB(t)+ \pi(t) \, dt],
\end{equation}
with
\[
\pi(t)=\frac{{\tr} \, (L\ga(t)+\ga(t) L^*)}{{\tr}\, \ga(t)}.
\]

As in the proof of Theorem \ref{LindStochLin1} above, we can
expand $\rho_0=\ga_0$ in a series
\[
 \ga_0=\sum_{k=1}^{\infty} p_k e_k\otimes \bar e_k
 \]
 with a non-negative sequence $\{p_k\}$ summing up to one
 and an orthonormal basis $\{e_k\}$.
Hence we can represent $\ga(t)$ as the convergence series of pure states
 \[
 \ga(t)=\sum_{k=1}^{\infty} p_k e_k(t)\otimes \bar e_k(t),
 \]
 with $e_k(t)$ solving the linear filtering equation for
pure states \eqref{eqqufiBlins}:
\begin{equation}
\label{eqinfdim}
de_k(t)=(-iH e_k(t)-\frac12 L^*L e_k(t))\,dt +Le_k(t) [dB(t)+\pi(t)\, dt].
\end{equation}
Here
\[
\pi(t)=\frac{\sum_{k=1}^{\infty} p_k (e_k(t), (L+L^*) e_k(t))}{\sum_{k=1}^{\infty} p_k \|e_k(t)\|^2}.
\]

It is a key observation that the infinite-dimensional system of SDEs \eqref{eqinfdim}
can be considered as a single SDE
with values in the  Hilbert space $l^2_{\HC}(\{p_k\})$ consisting of infinite sequences
$\e=(e_1, e_2, \cdots )$ of vectors from $\HC$ and equipped
with the norm
\[
\|\e\|^2=\sum_{k=1}^{\infty} p_k (e_k, e_k).
\]
Bounded operators in $\HC$ extend naturally (acting identically on each coordinate)
to bounded operators in $l^2_{\HC}(\{p_k\})$ with the preservation of norm.
In this notation system \eqref{eqinfdim} writes down as the SDE
\begin{equation}
\label{eqinfdim1}
d \e(t)=(-iH \e(t)-\frac12 L^*L \e(t))\, dt +L\e(t) \left[dB(t)+\frac{(\e, (L+L^*) \e)}{(\e, \e)} \, dt\right].
\end{equation}

This equation is the same as \eqref{eqqufiBlinsB} (though written in an enhanced Hilbert space). Hence
the coefficients of this equation are globally Lipschitz due to Lemma \ref{lemboundder}.
Therefore, it has the unique solution. And consequently, \eqref{Lindstochnormlinmix} has a unique solution $\ga(t)$,
and hence  $\rho(t)=\ga(t)/T(t)$ is the unique solution to the Cauchy problem for \eqref{Lindstochnorm1}.
\end{proof}

The full well-posedness of a problem includes also a statement on a continuous dependence of
the solution on initial data and parameters of the problem. We prove here the continuous dependence
on the Hamiltonian, which would be crucial for the next Section.

 \begin{theorem}
\label{LindStochnonLin2}
Under the assumption of Theorem \ref{LindStochnonLin1}
let us consider the Cauchy problem for equations
\begin{equation}
\label{Lindstochnorm11}
d\rho(t)=-i[H+H_j(t),\rho(t)]\, dt+\LC_L \rho (t)\, dt
+[L\rho(t)+\rho(t) L^*-\rho(t)\, {\tr} \, (L\rho(t)+\rho(t) L^*) ] dB(t),
\end{equation}
$j=1,2$, where $H_1(t), H_2(t)$ are continuous families of bounded self-adjoint
 operators in $\HC$.
Then, for the solutions $\rho_j(t)$, $j=1,2$, of these equations with the same
positive initial data $\rho_0$ of unit trace, one has the following estimate:
\begin{equation}
\label{Lindstochnorm12}
\E \|\rho_1(t)-\rho_2(t)\|_{H^{1,2}_s} \le \sqrt t C(t) \sup_{s\in [0,t]}\|H_1(s)-H_2(s)\|,
\end{equation}
where $C(t)$ is an increasing continuous function depending on $\|L\|$,
and where $\HC^{1,2}_s$ means of course either $\HC^1_s$ or $\HC^2_s$.
\end{theorem}

\begin{proof} By changing to the "interaction picture", that is, to equations
on the variable $e^{-itH}\rho(t) e^{itH}$ we can reduce the discussion to the case
of vanishing $H$. Time-dependence of $H_j$ and $L$ arising from this change does
not affect the argument. Therefore, without loss of generality we can set $H=0$.

Making the transformation to the equations in $l^2_{\HC}(\{p_k\})$,
as in the proof of the previous theorem,
 we can rewrite equations \eqref{Lindstochnorm11} as the equations
\begin{equation}
\label{eqinfdim11}
d \e^j(t)=(-iH_j \e^j(t)-\frac12 L^*L \e^j(t))\, dt +L\e^j(t)
\left[dB(t)+\frac{(\e^j, (L+L^*) \e^j)}{(\e^j, \e^j)} \, dt\right],
\end{equation}
with $j=1,2$.

As shown in the previous theorem, these equations are SDEs in a Hilbert space
with globally Lipschitz coefficients. Moreover, these two equations differ by
bounded linear terms. Hence it is a standard procedure
(see e.g. Proposition 7.1 in \cite{KolQuantLLN}) to derive an
estimate for $\E \|(\e^1-\e^2)(t)\|^2$ in terms of $\|H_1-H_2\|$.
We skip details here referring to \cite{Kol25a}, where they are given in full.
\end{proof}

\subsection{Semigroups generated by nonlinear filtering equations}

Let us look at the analytic properties of the semigroup $\Phi_tf(\rho)=\E_{\rho} f(\rho(t))$,
where $\rho(t)$ are solutions to the Cauchy problem of SDE \eqref{Lindstochnorm1}. As follows
from Theorem \ref{LindStochnonLin}, the semigroup $\Phi_t$ is a $C$-Feller semigroup in $C(S(\HC))$
(a semigroup of positivity preserving contractions). For finite-dimensional $\HC$ it implies directly
that smooth functions (of any order) represent invariant cores for generator \eqref{eqdifgener}.
However, for infinite-dimensional $\HC$, when searching for an invariant core
(where generator \eqref{eqdifgener} is defined), we face the same problem
as for the case of countable observation above, namely that smooth functions
 are not dense in $C(S(\HC))$. To tackle this issue, we shall reduce the action
 of $\Phi_t$ to the subspace $\hat C(S(\HC))\subset C(S(\HC))$ consisting of functions that
 are continuous in the topology of $\HC^2_s$. In other words, functions from $\hat C(S(\HC))$
 are restrictions to $S(\HC)$ of the functions from the space $C(\HC^2_sP^+)$ of continuous functions
 on the projective space $\HC^2_sP^+$ of lines in $\HC^2_s$ going through nonnegative operators. Therefore, yet equivalently,
 the space $\hat C(S(\HC))$ is isomorphic to the space $C(S^+(\HC^2_s))$ of continuous functions on
 the positive part $S^+(\HC^2_s)$ of the unit sphere in $\HC^2_s$.

It is an easy exercise to derive the evolution of states normalised in the sense of $\HC^2_s$
(not $\HC^1_s$, as we did earlier) from evolution \eqref{Lindstoch}. Namely,
from \eqref{Lindstochquadtr} we derive that
\[
d\|\ga\|_2^{-1}=-\frac{1}{\|\ga\|_2^3}(L+L^*,\ga^2)\, dY(t)
\]
\[
-\frac{1}{2\|\ga\|_2^3} \,{\tr}\, (2\ga L\ga L^*+L \ga L \ga + \ga L^* \ga L^*) \, dt
 +\frac32 \frac{1}{\|\ga\|_2^5} (L+L^*,\ga^2)^2\, dt,
 \]
 where we use scalar product in the space $\HC^2_s$, so that $(L+L^*,\ga^2)={\tr}((L+L^*)\ga^2)$.
Therefore for $\hat \ga=\ga/\|\ga\|_2$ we obtain
\[
d \, \hat  \ga=(-i[H,\hat \ga(t)] \, dt +\LC_L \hat \ga(t)) \, dt
\]
\[
-\frac12 \hat \ga \, {\tr} \,(2\hat \ga L\hat \ga L^*+ L \hat \ga L \hat \ga + \hat \ga L^* \hat \ga L^*) \, dt
\]
\[
-(L+L^*, \hat \ga^2) (L\hat \ga +\hat \ga L^*) \, dt+\frac32 \hat \ga \, (L+L^*, \hat \ga^2) ^2 \, dt
\]
\[
+[L\hat \ga(t)+\hat \ga(t) L^* -\hat \ga (L+L^*, \hat \ga^2)]\, dY(t).
\]

From this one derives
\[
d(\hat \ga, \hat \ga)=2(1-(\hat \ga, \hat \ga)) (L+L^*, \hat \ga^2)\, dY(t)
+2( \hat \ga , \LC_L \hat \ga(t)) \, dt
\]
\[
-(\hat \ga, \hat \ga) \, {\tr} \,(2\hat \ga L\hat \ga L^*+ L \hat \ga L \hat \ga + \hat \ga L^* \hat \ga L^*) \, dt
\]
\[
-2(L+L^*, \hat \ga^2)^2 \, dt +3 (\hat \ga, \hat \ga) \, (L+L^*, \hat \ga^2) ^2 \, dt
\]
\[
+ (L\hat \ga(t)+\hat \ga(t) L^* -\hat \ga (L+L^*, \hat \ga^2),L\hat \ga(t)+\hat \ga(t) L^* -\hat \ga (L+L^*, \hat \ga^2)) dt.
\]
\[
=(1-(\hat \ga, \hat \ga))\bigl[[2 (L+L^*, \hat \ga^2)\, dY(t)-4  (L+L^*, \hat \ga^2)^2\, dt
\]
\[
+ \, {\tr} \,(2\hat \ga L\hat \ga L^*+ L \hat \ga L \hat \ga + \hat \ga L^* \hat \ga L^*) \, dt\bigr].
\]
This preserves $(\hat \ga, \hat \ga)$, when it equals to one.

If we modify the above equation for $\hat \ga $ to (normalising $\ga$ in the nonlinear terms)
 \[
d \, \hat  \ga=(-i[H,\hat \ga(t)] \, dt +\LC_L \hat \ga(t)) \, dt
\]
\[
-\frac12 \hat \ga \, {\tr} \,(2\hat \ga L\hat \ga L^*+ L \hat \ga L \hat \ga + \hat \ga L^* \hat \ga L^*)
 \frac{1}{\|\hat \ga\|^2_2}\, dt
\]
\[
-\left(L+L^*, \frac{\hat \ga^2}{{\|\hat \ga\|^2_2}}\right) (L\hat \ga +\hat \ga L^*) \, dt
+\frac32 \hat \ga \, \left(L+L^*, \frac{\hat \ga^2}{{\|\hat \ga\|^2_2}}\right) ^2 \, dt
\]
\begin{equation}
\label{eqquantfilonhilbertsphere}
+\left[L\hat \ga(t)+\hat \ga(t) L^* -\hat \ga \left(L+L^*, \frac{\hat \ga^2}{{\|\hat \ga\|^2_2}}\right)\right]\, dY(t),
\end{equation}
then we get $d(\hat \ga, \hat \ga)=0$,
and the conservation of the norm becomes universal, but the evolution of operators
with unit Hilbert-Schmidt norm would not be changed.
Moreover, by Lemma \ref{lemboundderden},
equation \eqref{eqquantfilonhilbertsphere} has smooth coefficients
in $S^+(\HC)$. It follows that evolution given by \eqref{eqquantfilonhilbertsphere} is well-defined
and preserves smooth functions (of arbitrary order) on $S^+(\HC)$. Since $S^+(\HC)$
is a Hilbert-based infinite-dimensional manifold, the classes of smooth functions
$C^k(S^+(\HC))=C^k(\HC^2_sP^+)$ are dense
in the space of continuous functions for any $k$. But evolution  \eqref{Lindstochnorm1} is just the restriction
of evolution \eqref{eqquantfilonhilbertsphere} on $S^+(\HC)$, or equivalently on $\HC^2_sP^+$.
Identifying, as noted above, the spaces
 $\hat C(S(\HC))$, $C(\HC^2_sP^+)$ and $C(S^+(\HC^2_s))$ we get dense classes of invariant smooth
 (in the topology of $\HC^2_s$, of course) functions on $S(\HC)$.

Thus we have proved the following result.

\begin{theorem}
\label{nonLinsdesem}
The semigroup $\Phi_t$ is strongly continuous in the space $\hat C(S(\HC))$, having the spaces
$C^k(S^+(\HC))=C^k(\HC^2_sP^+)$ as invariant cores for any $k\ge 2$.
\end{theorem}

Since generator \eqref{eqmixgenerun} of quantum filtering process for mixed observation (and unitary coupling
operators with counting measurements) is seen to be obtained from \eqref{eqdifgener} by adding a bounded perturbing
term, the standard perturbation theory allows one to directly extend the previous theorem to the case of operators
\eqref{eqmixgenerun}  yielding the following result.

\begin{theorem}
\label{nonLinsdesemmix}
The semigroup $\Phi_t^{mix}$ generated by operator \eqref{eqmixgenerun}
is strongly continuous in the space $\hat C(S(\HC))$, having the spaces
$C^k(S^+(\HC))=C^k(\HC^2_sP^+)$ as invariant cores for any $k\ge 2$.
\end{theorem}

\subsection{Classical Hamiltonians}

When $H$ is unbounded, one can work with the interaction representation
 or with mild forms of filtering equations,
as we did above, reducing the analysis to bounded operators. However, it is of course desirable to solve
the corresponding filtering equations in their original form. For this to be possible, one has to ensure that
the evolution remains inside the domain of $H$ for all times.

For completeness, we briefly present here the simplest basic examples (which will be used
also when dealing with the propagation of chaos) referring to \cite{Kol25b}
for the full story including unbounded $L$.

The classical Hamiltonian of quantum mechanics acts in $L^2(\R^d)$ as the operator
\[
Hf(x)=-\frac12 \De f(x) +V(x) f(x),
\]
where $V(x)$ is a function in $\R^d$, which is also identified (as usual) with the operator
of multiplication by this function.
Recall that a standard example of the condition ensuring that $H$ is self-adjoint
in $L^2(\R^d)$ on the domain being the Sobolev space $W^2(\R^d)$ is that $V$ is bounded
for $d\le 2$ and  that $V=V_1+V_2$,
where $V_2$ is bounded and $V_1\in L^p(R^d)$ for $d\ge 3$ and $p>d/2$, see e.g. \cite{ReedSimon}
and more general situations in \cite{Yajima}.

\begin{theorem}
\label{thclassichamfilt}
Let $V$ be a measurable function on $\R^d$ such that the operator $H=-\frac12 \De +V(x)$ is self-adjoint
in $L^2(\R^d)$ on the domain being the Sobolev space $W^2(\R^d)$. Let $L$ be a multiplication operator
on the real function $L$ that belongs to $C^2(\R^d)\cap W^2(\R^d)$. For dimensions $d=1,2,3$ it
is sufficient to assume $L\in C^1(\R^d)\cap W^2(\R^d)$. Then the equations
\eqref{eqqufiBlins} and \eqref{eqqufiBnonlinsn} are well-posed in the sense that for any initial condition
from $W^2(\R^d)$ there exists a unique global solution of the corresponding Cauchy problem that belongs
$W^2(\R^d)$ for all times.
\end{theorem}

\begin{proof}
The simplest proof is via the mild forms
\eqref{eqqufiBlinsm} and \eqref{eqqufiBnonlinsm}. Though the well-posedness
was proved via the interaction representation, we note that it can be also obtained directly from
these mild forms via the standard Banach fixed-point principle. Moreover,
from the condition of the theorem it follows that the semigroup $e^{iHt}$
is also bounded and strongly continuous in $W^2(\R^d)$ (see e.g. \cite{Yajima}).
Moreover, the operators of multiplication on $L$ and $L^2$ are bounded in $W^2(\R^d)$.
The same holds under simplified assumptions for $d=1,2,3$, because in these dimensions
$W^2(\R^d)\in L^{\infty}(\R^d)$ (see e.g. \cite{ReedSimon}).  Therefore,
the same Banach fixed-point principle applied in $W^2(\R^d)$ allows one
to conclude that if the initial condition was in $W^2(\R^d)$, then it will stay there for all times.
Consequently, the integrands on the r.h.s. of equations \eqref{eqqufiBlinsm} and \eqref{eqqufiBnonlinsm}
become differentiable in $t$, and consequently, solutions to these equations become also solutions to
\eqref{eqqufiBlins} and \eqref{eqqufiBnonlinsn}.
\end{proof}

Analogous results hold for the equations on mixed states. Recall that we proved
the well-posedness of the interaction representation (and hence also mild form)
of equation  \eqref{Lindstoch} in the Hilbert space of self-adjoint
Hilbert Schmidt operators $\HC^2_s$. In the setting of Theorem \ref{thclassichamfilt} this space
coincides with the Hilbert subspace  $L^2_s((\R^d)^2)$ of symmetric functions in
$L^2((\R^d)^2)$ (integral kernels such that $\ga(x,y)=\overline{\ga(y,x)}$).
Moreover, the domain of the generator $-i[H, \ga]$ of the semigroup of operators
$\ga \mapsto e^{-iHt}\ga e^{iH t}$ in $\HC^2_s$  is given by the Sobolev space $W^2((\R^d)^2)$.
Consequently, the same argument as above gives the following.

\begin{theorem}
\label{thclassichamfiltmix}
Under the assumptions of Theorem \ref{thclassichamfilt},  equations
  \eqref{Lindstoch} and \eqref{Lindstochnorm1} are well posed in $\HC_s^2$ in the sense that for any
  $\ga_0\in W^2(\R^{2d})$ (resp. $\ga_0\in W^2(\R^{2d})\cap S(L^2(\R^d))$ there
 exist unique global solutions of the corresponding Cauchy problems that belong
$W^2((\R^d)^2)$ for all times.
\end{theorem}

Similar results hold for the filtering equations of counting observations.


\subsection{Appendix: some trace inequalities}

\begin{prop}
If $A$ is a self-adjoint Hilbert-Schmidt operator and $B$ a bounded operator, then
\begin{equation}
\label{eqmyineqtr}
2|{\tr} \, (ABAB^*)|\le {\tr} \, [A^2(BB^*+B^*B)],
\end{equation}
and
\begin{equation}
\label{eqmyineqtr1}
|{\tr} \, (AB AB +AB^*AB^*)|\le {\tr} \, [A^2 (BB^*+B^*B)].
\end{equation}
\end{prop}

\begin{proof}
By approximation it is reduced to finite-dimensional situation. The diagonalization procedure reduces
the problem to the case when $A$ is a diagonal matrix with real numbers $a_i$ on the diagonal. Then
 \[
 2{\tr} \, (ABAB^*)
=2\sum a_ib_{ij}a_j\bar b_{ij}=\sum a_i a_j (|b_{ij}|^2+|b_{ji}|^2)
\]
\[
=2\sum_i a_i^2 |b_{ii}|^2+2\sum_{i<j}a_ia_j(|b_{ij}|^2+|b_{ji}|^2).
\]
The r.h.s. of \eqref{eqmyineqtr} equals
\[
\sum a_i^2 (|b_{ij}|^2+|b_{ji}|^2)
=2\sum_i a_i^2 |b_{ii}|^2+ \sum_{i < j}(a_i^2+a_j^2)(|b_{ij}|^2+|b_{ji}|^2).
\]
Thus \eqref{eqmyineqtr} holds, because $2|a_ia_j| \le a_i^2 +a_j^2$.

Inequality \eqref{eqmyineqtr1} rewrites as
\[
2\sum_i a_i^2 |Re (b_{ii}^2)|+4\sum_{i<j}|a_ia_j| \, |Re (b_{ij} b_{ji})|
\le 2\sum_i a_i^2 |b_{ii}|^2+ \sum_{i < j}(a_i^2+a_j^2)(|b_{ij}|^2+|b_{ji}|^2),
\]
which easily seen to hold.
\end{proof}

In particular, for self-adjoint $B$ it follows that

\begin{equation}
\label{eqmyineqtr2}
|{\tr} \, (ABAB)|\le {\tr} \, (A^2 B^2).
\end{equation}


\section{Derivation}
\label{secder}


As we noted, the theory of quantum filtering was built initially in papers
\cite{Bel87}, \cite{Bel88}, \cite{Bel92}
on the basis of quantum stochastic calculus.
A well written review of this development is given in \cite{BoutHanJamQuantFilt}.
Afterwards, there appeared alternative approaches to the derivations of the filtering equations,
see \cite{Attal}, \cite{BelKol},
\cite{Pellegrini}, \cite{BarchBel}, \cite{Holevo91}, \cite{KolQuantFrac} and references therein.
Some further extensions can be found in \cite{BarndLoub}, \cite{Loubenets}.
The most straightforward (and intuitive) derivation seems to be via the limit
of successive instantaneous measurements,
when the times between these measurements tend to zero. The complete rigorous derivation
of the filtering equations via this method  in both diffusive and counting observations
was performed for finite-dimensional quantum mechanics in \cite{Pellegrini} and
\cite{Pellegrini10}. In \cite{KolQuantFrac} the author added the explicit rates of convergence
and indicated the ways to infinite-dimensional extensions.
Using the well-posedness results above, we present here, seemingly for the first
time, the complete rigorous derivation of the filtering equations for the general
infinite-dimensional case. For completeness,
we also discuss by-passing the famous quantum Zeno paradox.


\subsection{Quantum measurement}

If two quantum systems living in Hilbert spaces $\HC$ and $\HC_1$
are brought to interaction, the combined system
has the tensor product Hilbert space  $\HC \otimes \HC_1$ as the state space.
We will be interested now only in the case,
when the Hilbert space $\HC_1$ is finite-dimensional, that is, $\HC_1=\C^m$.
In this case $\HC\otimes \HC_1$ can be identified with the Hilbert-space $\HC^m$
of vectors $(h_0, h_1, \cdots, h_{m-1})$ with
coordinates $h_j\in \HC$, so that if $h\in \HC$ and $f=(f_0, \cdots, f_{m-1})\in \C^m$, then
$h\otimes f=(f_0 h, \cdots, f_{m-1} h)$.
Similarly, operators $A$ in $\HC \otimes \HC_1$ are
given by matrices $\{A^i_j\}$, $i,j=0, \cdots , {m-1}$,
with each entry $A^i_j$ being an operator in $\HC$.
A product $A\otimes B$ of two operators $A$ and $B$ acting in $\HC$ and $\HC_1$ respectively
is defined generally by its action on the tensor products as
\[
(A\otimes B)(e\otimes f)=Ae\otimes Bf.
\]
In case $\HC=\C^m$, it means that if $B$ is given by an $m\times m$-matrix with entries $B_{ij}$,
then $A\otimes B$ is given by the matrix with the entries $AB_{ij}$.

An operator $A$ in $\HC$ has the natural lifting $A\otimes I$ (where $I$ is the unit operator) to $\HC\otimes \HC_1$.
Similarly an operator $B$ in $\HC_1$ has the natural lifting $I\otimes B$ to $\HC\otimes \HC_1$.

The key notion of the theory of interacting systems is that of the {\it partial trace}. For an operator
$A$ in $\HC \otimes \HC_1$ given by its matrix $\{A^k_j\}$, the partial trace with respect to $\HC_1$
is the operator in $\HC$ defined as
 \begin{equation}
 \label{eqdefparttr}
 {\tr}_{p1} A=\sum_k A^k_k.
 \end{equation}
This  partial trace is interpreted as the state of the first system given the state of the coupled one.
 Similarly, the partial trace with respect to $\HC$ is the operator
 in $\HC_1$ defined as
 \[
 ({\tr}_{p0} A)^i_j={\tr} A^k_j.
 \]
 (of course, whenever all these traces exist).
 Clearly,
 \[
 {\tr} ({\tr}_{p0} A)= {\tr} ({\tr}_{p1} A)= {\tr} (A).
 \]

{\it Physical observables} in the Hilbert space of a quantum system are
given by self-adjoint operators $A$ in it. If $A$ has a discrete spectrum,
then $A$ has the spectral decomposition $A=\sum_j \la_j P_j$, where
 $P_j$ are orthogonal projections on the eigenspaces of $A$ corresponding to
 the eigenvalues $\la_j$. According to the {\it basic postulate of quantum measurement}
 \index{basic postulate of quantum measurement}, measuring observable $A$ in a state $\ga$
(often referred to as the {\it Stern-Gerlach experiment}\index{ Stern-Gerlach experiment})
can yield each of the eigenvalue $\la_j$ with the probability
   \begin{equation}
\label{eqantiunitdress1}
{\tr} \, (\ga P_j)={\tr} \, (P_j \ga P_j),
\end{equation}
 and, if the value $\la_j$ was obtained, the state
of the system changes (instantaneously) to the reduced state
\[
P_j\ga P_j/ {\tr} \, (\ga P_j).
\]
In particular, if the state $\rho$ was pure, $\ga=|\psi\rangle \langle \psi|$, then the
probability to get $\la_j$ as the result of the measurement becomes $(\psi_,P_j\psi)$
and the reduced state also remains pure and is given by the vector $P_j\psi$.
%If the interaction with the apparatus was preformed 'without reading the results',
%the state $\rho$ is said to be subject to a {\it non-selective measurement}
%\index{non-selective measurement} that changes $\ga$ to the state $\sum_j P_j\rho P_j$.

It is physically more natural to make indirect measurement of a chosen quantum system,
which we sometimes refer to as an atom, where this system is coupled (brought to interaction) with
some measuring apparatus, whose pointer positions (states) we can read.  More precisely,
{\it indirect measurements} of an atom in the initial space $\HC$,
are organised as follows.
One couples the atom with another quantum system, a measuring devise,
also referred to an ancilla system or s reservoir, specified
by another Hilbert space $\HC_1$, which in this paper we choose to be finite-dimensional: $\HC_1=\C^m$.
 Namely the combined system lives in the tensor product
 Hilbert space $\HC\times \HC_1$ and its evolution is given by  certain self-adjoint
 operator $H$ in  $\HC\times \HC_1$. In the measuring device some fixed vector $\varphi \in \HC_1$
 is chosen, called the vacuum and interpreted as the stationary state of the devise when
 no interaction is involved. The corresponding density matrix will be denoted $\Om=|\varphi \rangle \langle \varphi|$.
 Indirect measurements of the states of the atom
 are performed by measuring the coupled system via an observable of the second system
 and then projecting the resulting state to the atom via the partial trace.

Namely, it is given by an operator $R$ in $\HC_1$ with the spectral decomposition
$R=\sum_j \la_j P_j$, where $P_j=f_j\otimes \bar f_j$ and $f_j$ the normalised eigenvectors of $R$,
 and it is performed in two steps: given a mixed state $\ga$ in  $\HC\times \HC_1$
one performs a measurement of $R$ lifted as $I\otimes R$ to $\HC\times \HC_1$
yielding values $\la_j$ and new states
\[
(I\otimes P_j)\ga (I\otimes P_j)/ {\tr} \, (\ga  (I\otimes P_j))=(f_j, \ga f_j) \otimes P_j
\]
with the probabilities
\[
p_j= {\tr} \, (\ga  (I\otimes P_j))={\tr} \, (f_j, \ga f_j),
\]
and then one projects these states to $\HC$
via the partial trace producing the states
\begin{equation}
\label{eqindirmeas}
\frac{{\tr}_{p1} [(I\otimes P_j)\ga (I\otimes P_j)]}{ {\tr} \, (\ga  (I\otimes P_j))}
=\frac{(f_j, \ga f_j) \otimes P_j}{{\tr}\, (f_j, \ga f_j)},
\end{equation}
where, for an operator-valued matrix $\ga$ and a vector $f=\{f^l\}\in \C^m$ we denoted
\[
(f,\ga f)=\sum_{k,l} \bar f^k\ga_{kl} f^l.
\]

\subsection{Quantum Zeno paradox and watchdog effect}

A natural idea of organising continuous measurement is via a limit of sequential instantaneous
measurements with times between these measurements tending to zero. Let us see what comes out of it.
Suppose a starting point $\psi_0$ of our atom in $\HC$ is chosen, between measurements
 the atom evolves according to the Hamiltonian $H$ and each measurement is made according
 to the operator $P_0$, which is the projection on $\psi_0$. In other words,
 we check every time whether our state is $\psi_0$ or is orthogonal to it.
 Assume $H$ is bounded. Then starting with $\psi_0$ we move in a small time $t$ to
 \[
 e^{iHt}\psi_0\sim \frac{\psi_0+iHt \psi_0}{\|\psi_0+iHt \psi_0\|}\sim (\psi_0+iHt \psi_0)(1+O(t^2)).
 \]
 Therefore the probability to return to $\psi_0$ after the first measurement at time $t$ is
 \[
 |(e^{iHt}\psi,\psi_0)|^2(1+O(t^2))\sim |(\psi_0+iHt \psi_0,\psi_0)|^2(1+O(t^2))=1+O(t^2).
 \]
 Fix now a time interval of length $T$, choose natural $n$,  and let us make instantaneous
 measurements at times $t_k=(k/n)T$. Then the probability to remain in $\psi_0$ after all $n$
 measurements is $(1+O(k^2T^2/n^2))^n$, which clearly tends to $1$, as $n\to \infty$.
 This is the {\it quantum Zeno paradox} -- when continuous measurement is performed in this way,
 the atom just remains in its initial state independently of the law of its evolution.
 This paradox was suggested in \cite{Zeno}, where we refer to for the corresponding
  argument in case of unbounded (in fact, semibounded) Hamiltonian operator.
 By obvious reasons this effect was also named {\it watch-dog effect}, or yet alternatively
 expressed by  wording "watched kettle never boils".

 To avoid this effect, one has to scale appropriately the dynamics in order to
 make the effect of evolution comparable with the effect of measurements.
 We show below how this is done in the framework of indirect measurements.

 Let us mention however that Zeno paradox above was linked to the assumptions that
 we keep measuring the same projection operator $P_0$ all the times. An alternative way to deviate
 from Zeno paradox is to smoothly change the measuring observable from $P_0$ to $P_t=U(t)P_0U^*(t)$
 with a smooth family of unitary operators $U(t)$. Then similar calculations will show that
 in the limit of continuous measurement the system will move through the eigenstates of $P_t$
 with eigenvalue $1$. This effect was discovered in \cite{Anti-Zeno} and called
 {\it Anti-Zeno paradox}.

\subsection{Markov chains of sequential indirect observations}
\label{secMarchainindirodserv}

We describe now the
Markov chains of sequential indirect observations, rather standard by now,
in discrete and continuous time.


The discrete time {\it Markov chain of successive indirect observations} (or measurements) evolves according
to the following procedure specified by a triple: a self-adjoint  operator $A$ in $\HC\times \HC_1$,
a self-adjoint operator $R$ in $\HC_1$ and the vacuum vector $\Om$ in $\HC_1$.
(i) Starting with an initial state $\rho$ of $\HC$ one couples it with the device in its vacuum state $\Om$
producing the state $\ga=\rho\otimes \Om$ in $\HC\times \HC_1$,
(ii) During a fixed period of time $h$ one evolves the system according to the operator $A$ producing the state
$\ga_h=e^{-ihA}\ga e^{ihA}$ in $\HC\times \HC_1$,
(iii) One performs the indirect measurement with the state $\ga_h$ yielding the states
\begin{equation}
\label{eqMarkchain}
\rho_h^j=\tilde \rho_h^j/ p_j(h)
\end{equation}
where
\[
\tilde \rho_h^j ={\tr}_{p1} (I\otimes P_j)\ga_h (I\otimes P_j)
\]
\begin{equation}
\label{eqMarkchain0}
={\tr}_{p1} (I\otimes P_j)e^{-ihA}(\rho\otimes \Om) e^{ihA} (I\otimes P_j)
=(f_j, e^{-ihA}(\rho\otimes \Om) e^{ihA} f_j) \otimes P_j,
\end{equation}

are the corresponding non-normalised states, and
\begin{equation}
\label{eqMarkchain1}
p_j(h)={\tr} \, \tilde \rho_h^j
={\tr} \, (e^{-ihA}(\rho\otimes \Om) e^{ihA}  (I\otimes P_j))
={\tr}\, (f_j, e^{-ihA}(\rho\otimes \Om) e^{ihA} f_j)
\end{equation}
are the probabilities of obtaining these states as the result of the measurement.

Then the same repeats starting with $\rho_h$ as the initial state.  Let us denote $U_h$ the transition
operator of this Markov chain that acts on the set of continuous functions on $S(\HC)$ as
\begin{equation}
\label{eqMarkchain2}
U_h f(\rho)=\E f(\rho_h)=\sum_j p_j(h) f(\rho_h^j).
\end{equation}
Similarly one can define the continuous time
{\it Markov chain of successive indirect observations} (or measurements)
$O^{\rho}_{t,h}$ and the corresponding Markov semigroup $T_t^h$ on $C(S(\HC)$
evolving according to the same rules, with only difference that the times between successive measurements
are not fixed, but represent independent exponential random  variables $\tau$ with
the intensity $1/h$: $\P(\tau>t)=e^{-t/h}$.
The generator $L^h$ of this Markov process is bounded in $C(S(\HC))$ and acts as
\begin{equation}
\label{eqMarkchain3}
L^hf(\rho)=\frac{(U_hf-f)(\rho)}{h}=\frac{1}{h}\sum_j p_j(h) (f(\rho_h^j)-f(\rho)).
\end{equation}

All "quantum content" of the theory is now captured in the explicit formula \eqref{eqMarkchain}.
What follows will be
the pure classical probability analysis of these Markov chains, their scaling limits and control.

\begin{remark}
\label{redtimedepder}
Let us comment on the straightforward extension of the scheme above to the case of time dependent family
of operators $A$. Namely, suppose that instead of $A$ we have a family of self-adjoint
operators $A(t)$ in $\HC\otimes \HC_1$. Then in the discrete scheme we perform measurements at times
 $t_k=kh$ and between times $t_k$ and $t_{k+1}$ the system evolves according to $A(t_k)$ producing the state
$\ga_h=e^{-ihA(t_k)}\ga e^{ihA(t_k)}$ in $\HC\times \HC_1$. The rest is the same.
Of course, the corresponding  operators \eqref{eqMarkchain2} and \eqref{eqMarkchain3}
become time-dependent. We mention this extension, because the reduction
of the case of unbounded Hamiltonian via the interaction representation leads
necessarily to time-dependent families $A(t)$.
\end{remark}

In what follows, as the vacuum vector we shall choose the first basis vector $e_0=(1,0,\cdots , 0)$
 of our measuring devise Hilbert space $\HC=\C^m$, so that $\Om =e_0\otimes e_0$
 is the orthogonal projection to $e_0$.
Moreover, to shorten formulas we shall further on write down $m\times m$-matrices
 $D$ representing operators in $\HC_0\otimes \C^m$
 in the concise $2\times 2$-form, that is, as
 \[
R= \begin{pmatrix}
A & C   \\
D & B
\end{pmatrix},
\]
where $B$ is an operator-valued  $(m-1)\times (m-1)$-matrix, $C$ and $D$ are a row
and a column operator-valued vectors respectively.
 For instance, for an operator $D$ in $\HC_0$, we shall write
 \[
D\otimes \Omega
=\begin{pmatrix}
D & 0   \\
0 & 0
\end{pmatrix}.
\]

As all sequential measurements go through projection to the vacuum, the parts of a Hamiltonian operator $H$
in $\HC_0\otimes \C^m$ (describing the dynamics) that do not involve the vacuum are irrelevant for the limiting
dynamics (that we are interested in). Therefore, for the purpose of deriving the limiting dynamics
of continuous observation, we can and will consider only  Hamiltonians
with matrices of the form
\[
A= \begin{pmatrix}
H & -iL_1^* & -iL_2^* &\cdots & -iL_{m-1}^*  \\
iL_1 & 0 & 0 & \cdots & 0 \\
&  &  \cdots &  & \\
iL_{m-1} & 0 & 0 & \cdots & 0
\end{pmatrix},
\]
or better in our concise notations,
\[
A= \begin{pmatrix}
H & -iL^*   \\
iL & 0
\end{pmatrix}.
\]
Here $H$ is a self-adjoint operator in $\HC$ describing
the dynamics of our atom without observation and the vector-valued operator $L=(L_1, \cdots , L_{m-1})$
in $\HC$ describes the interaction of the atom with the measuring devise.

We are aiming at calculating the small time asymptotics of the
Markov transition operators defined by \eqref{eqMarkchain}.

We shall performed all calculations as if both $H$ and $L$ are bounded. However, as was mentioned earlier,
the case of unbounded $H$ reduces to the analysis of vanishing $H$ but with the time-dependent
$L(t)=e^{itH}Le^{-iHt}$ (see Remark \ref{redtimedepder}). Therefore, the only real restriction
in our derivation is the boundedness of $L$.

\begin{remark}
One can check that if one would write some arbitrary matrix $B$ instead
of the zero block in the expression for $H$, then
this $B$ would not contribute to the limiting evolution obtained below, see [Kol].
\end{remark}

The main idea for obtaining sensible asymptotic limits (thus avoiding Zeno paradox)
is to enhance the interaction part $C$ of the Hamiltonians
 replacing it by the scaled version $C/\sqrt h$. Thus we choose the Hamiltonian in the form
\begin{equation}
\label{interham}
A= \begin{pmatrix}
H & -iL^*/\sqrt h   \\
iL/\sqrt h & 0
\end{pmatrix}
\end{equation}

 \begin{remark}
 The necessity of scaling is clear from Zeno's paradox. In principle one can suggest
 scaling $L$ as $L/h^{\al}$ with some $\al>0$. However, as calculations would show, choosing $\al>1/2$ would force $L$
 to disappear from the limit (measurement would have no effect in the limit), and choosing $\al<1/2$
 would force $H$ to disappear from the limit,
  only with $\al=1/2$ a sensible contributions from both $H$ and $L$ remain.
 \end{remark}


In the calculations below we shall use the following simple general small time asymptotic
formula for the evolutions $e^{-itA}$:
\[
e^{-itA}\ga e^{itA}
=(1-it A -\frac12 t^2 A^2)\ga (1+it A-\frac12 t^2 A^2)+O(t^3 \|A\|^3)
\]
\[
=\ga-it [A,\ga]-\frac12 t^2 A^2\ga -\frac12 t^2 \ga A^2+t^2 A\ga A+O(t^3 \|A\|^3)
\]
\begin{equation}
\label{smalltimegroup}
=\ga-it [A,\ga]+t^2 (A\ga A-\frac12 \{A^2,\ga\})+O(t^3\|A\|^3).
\end{equation}

To use this formula, for a mixed state $\rho$ in $\HC$, we calculate its ingredients as follows:

\[
\left[A, \rho \otimes \Om \right]
=\begin{pmatrix} [H,\rho] & i\rho L^*/\sqrt h \\ iL\rho/\sqrt h & 0 \end{pmatrix},
\]
\[
A (\rho \otimes \Om) A
=\begin{pmatrix} H\rho H & -i H\rho L^*/\sqrt h\\ iL\rho H/\sqrt h & L\rho L^*/h \end{pmatrix},
\]
\[
A^2 =\begin{pmatrix} H^2+L^*L/h & -iH L^* /\sqrt h \\ iL H/\sqrt h & LL^*/h \end{pmatrix},
\]
\[
\{A^2, \rho \otimes \Om\}=\begin{pmatrix} \{H^2+L^*L/h,\rho\} & -i\rho H L^*/\sqrt h \\ iL H \rho/\sqrt h  & 0 \end{pmatrix}.
\]
Using \eqref{smalltimegroup},
%and keeping terms of order not exceeding $t$
we get the approximation
\begin{equation}
\label{eqdressedrho1}
e^{-ihA} (\rho\otimes \Om) e^{ihA}
=\begin{pmatrix} \rho -ih[H, \rho] -\frac12 h\{L^*L,\rho\} &  \sqrt h \rho L^* \\ \sqrt h L\rho & hL\rho L^* \end{pmatrix}
+O(h^{3/2}),
\end{equation}
which is the key formula for what follows.

Let us stress for clarity that in this concise formula it is meant
 that $L^*L=\sum_j L_j^*L_j$ and $L\rho L^*$ is the square matrix
with the entries $L_k \rho L_j^*$.

\subsection{Counting observation}
\label{seccountcase}

As it turns out (may be not intuitively quite clear why),
the limiting processes are quite different depending on whether
the set of eigenvectors of the operator $R$ (observable of the measuring devise) includes the vacuum or not.
Let us start with the first case, when $e_0$ is an eigenvector of $R$.
Then the unitary operator $U$ in $\C^m$ transferring the standard basis
of $\C^m$ to the basis of eigenvectors of $R$ can be chosen in the form

\begin{equation}
\label{eqdiagonproj}
U= \begin{pmatrix} 1 & 0 \\ 0 & u \end{pmatrix},
\end{equation}
with a unitary operator $u$ in $\C^{m-1}=e_0^{\perp}$ -- the orthogonal complement of the vacuum vector.

It is then seen that changing the standard basis of $\C^m$ to the new one via the transformation $U$ would
transfer $A$ to the operator of the same form with the same $H$, but with the vector $uL$ instead of $L$.
This means that in the case, when $e_0$ is an eigenvector of $R$, we can consider $R$ to be diagonal without
essential loss of generality, as more general case would result only in a linear
transformation of $C$. Therefore, let us assume that $R=\sum_j \la_j P_j$ is diagonal, so that each $P_j$
is the projection on the standard basis vector $e_j$.

 Then the non-normalized new states are
\[
\tilde \rho_0=(I\otimes P_0) e^{-ihH} (\rho \otimes \Om) e^{ihH} (I\otimes P_0)
= \rho -ih[H, \rho] -\frac12 h\{L^*L,\rho\},
\]
\[
\tilde \rho_j=(I\otimes P_j) e^{-ihH} (\rho \otimes\Om) e^{ihH} (I\otimes P_j)
= hL_j\rho L_j^*, \quad j>0,
\]
occurring with the probabilities
\[
p_0=1-h \, {\tr} \, (L^*L \rho), \quad p_j=h \, {\tr}\, (L_j^*L_j \rho), j>0.
\]
Thus equation \eqref{eqMarkchain2} becomes
\[
U_h f(\rho)=\E f(\rho_h)=(1-h \, {\tr} \, (L^*L \rho))f
\left( \frac{\rho -ih[H, \rho] -\frac12 h\{L^*L,\rho\}}{1-h \, {\tr} \, (L^*L \rho)}\right)
\]
\begin{equation}
\label{eqMarkchaindiag}
+h \sum_{j>0} {\tr} (L_j^*L_j \rho) f\left(\frac{L_j \rho L_j^*}{{\tr}\, (L_j^*L_j \rho)}\right).
\end{equation}

Aiming at using Proposition \ref{propconvsemigr} (ii) we are looking for
the limit of the operator $(U_h-1)/h$ for $h \to 0$.

Denoting $T_j=  {\tr}\, (L_j^*L_j \rho)$ for $j>0$ and $T=\sum_j T_j$ we obtain
\[
\frac{U_h-1}{h} f(\rho)=\frac{1}{h} (1-hT)
\left[f \left( \frac{\rho -ih[H, \rho] -\frac12 h\{L^*L,\rho\}}{1-h T}\right)-f(\rho)\right]
\]
\[
+\sum_{j>0} T_j \left[f(\frac{L_j\rho L_j^*}{T_j})-f(\rho)\right].
\]
Since
\[
\frac{\rho -ih[H, \rho] -\frac12 h\{L^*L,\rho\}}{1-h T}-\rho
=\frac{-ih[H, \rho] -\frac12 h\{L^*L,\rho\}+h T \rho}{1-h T},
\]
we get by the Taylor expansion that, up to terms of order $\sqrt h$,
\[
\frac{U_h-1}{h} f(\rho)\approx \AC_{count} f(\rho)
\]
with
\begin{equation}
\label{eqjumpgenerrep}
\AC_{count}f(\rho)=-(f'(\rho), i[H, \rho] +\frac12 \{L^*L,\rho\}-\rho T)
+\sum_j T_j \left[f(\frac{L_j\rho L_j^*}{T_j})-f(\rho)\right],
\end{equation}
which is exactly operator $\AC_{count}$ from \eqref{eqjumpgener}.

Summarising, we conclude the following.

\begin{lemma}
\label{lemmaongencount}
Under the setting considered,
\begin{equation}
\label{eqjumpgener0}
\frac{U_h-1}{h} f \to \AC_{count}f
\end{equation}
for $f\in C^1_{Gat*}(X)$ (see Theorem \ref{JumpWellPDE1} for this notation),
with $\AC_{count}$ given by \eqref{eqjumpgenerrep}. Moreover,
\begin{equation}
\label{eqjumpgener1}
\|\frac{U_h-1}{h} f -\AC_{count}f\|\le \sqrt h \ka \|f\|_{C^2(S(\HC_0))}
\end{equation}
for $f\in C^2(S(\HC_0))$ and a constant $\ka$.
\end{lemma}

We can now complete our derivation of quantum filtering equations for continuous observations.


\begin{theorem}
\label{thBeleqcount}

(i) The scaled discrete semigroups $(U_h)^{[t/h]}$ converge to the semigroup $T_t$,
generated by operator \eqref{eqjumpgener} or \eqref{eqjumpgenerrep}
(according to Theorems \ref{JumpWellPDE} and \ref{JumpWellPDE1})
as $h\to 0$, so that the corresponding
processes converge in distribution; the scaled semigroups $T_t^h$ generated
 \eqref{eqMarkchain3} converge to the semigroup $T_t$, as $h\to 0$,
so that the corresponding
processes converge in distribution.

(ii) The rates of convergence can be given for smooth functions:
\begin{equation}
\label{eq1thBeleqcount}
\|(U_h)^{[t/h]}f -T_tf\| \le \sqrt h t M e^{Mt} \|f\|_{C^2(S(\HC))},
\end{equation}
and
\begin{equation}
\label{eq2thBeleqcount}
\|T_t^h f -T_tf\| \le \sqrt h t Me^{Mt} \|f\|_{C^2(S(\HC))}.
\end{equation}
with a constant $M$ depending on the norms of $H$ and $L$.
\end{theorem}

\begin{proof}
(i) This is a consequence of Lemma \ref{lemmaongencount}, Theorem \ref{JumpWellPDE1}
 and the general result on the convergence
of semigroups in terms of the convergence of their generators, see e.g. Theorems 19.27 and  19.28 of \cite{Kallen}.

(ii) This is a consequence of Lemma \ref{lemmaongencount}, Theorem \ref{JumpWellPDE}, Proposition \ref{propconvsemigr},
and the observation that \eqref{eq5propconvsemigr}
holds here with  the triple of spaces $C^2(S(\HC))\subset C^1(S(\HC)) \subset C(S(\HC))$.
\end{proof}


In order for Part (ii) of the Theorem to have a content, it is necessary for the classes
of smooth functions  $C^2(S(\HC))$  to be dense in $C(S(\HC))$. As we noted already, this is not clear in general.
Thus the rates of convergence in (ii) really make sense only either for finite-dimensional $\HC$ or for the case
of unitary operator $L$ (as seen from Proposition \ref{JumpWellPDE2}).


\subsection{Diffusive and mixed observation}
\label{secdifcase}

Let us turn to the second case, when $e_0$ is not an eigenvector of $R$.
Then the number $k$ of eigenvectors of $R$ having a non-vanishing projection
on $e_0$ is not less than $2$. Ordering eigenvectors of $R$ in such a way that
these $k$ eigenvectors take the first $k$ places, we conclude that in the
coordinate representation of the eigenvectors of $R$
\[
r_j=\sum_{l=0}^{m-1} r_j^l e_l,
\]
we have $r_j^0\neq 0$ exactly for $j=0, \cdots, k-1$.
Just for simplicity of writing assume that all coordinates here are real.

Recalling \eqref{eqdressedrho1} we derive the
new non-normalised states
\[
\tilde \rho_j =(r_j, e^{-ihA} (\rho\otimes \Om) e^{ihA} r_j)
\]
\[
=(r_j^0)^2 (\rho -ih[H, \rho] -\frac12 h\{L^*L,\rho\})
+r^0_j \sqrt h (\rho \tilde L_j^*+\tilde L_j \rho)+h\tilde L_j \rho \tilde L_j^*
+O(h^{3/2})
\]
where
\[
\tilde L_j= \sum_{l=1}^{m-1} r^l_j L_l.
\]

Introducing, as above, the notations $T_j={\tr}\, (\tilde L_j^* \tilde L_j \rho)$ for $j>0$ and $T=\sum_j T_j$
and denoting $\Om_j=  {\tr} (\rho \tilde L_j^* + \tilde L_j\rho)$ and observing that
$T=\sum_j {\tr} \, L_j^* L_j\rho $ we obtain
the probabilities of the occurrence of these states (up to higher order in $h$):
\[
p_j={\tr} \, \tilde \rho_j=(r_j^0)^2  (1-hT)+\Om_j r^0_j \sqrt h  +h T_j,
\]

In particular, for $j\ge k$,
\[
\tilde \rho_j =h\tilde L_j \rho \tilde L_j^*, \quad p_j=hT_j, \quad \rho_j=\frac{\tilde L_j \rho \tilde L_j^*}{T_j}.
\]

Next, for arbitrary numbers $a,b,c$, one can write up to terms of order $t$, that
\[
\frac{1}{a+b\sqrt t +ct}
=\frac{1}{a} \frac{1}{1+(b/a) \sqrt t+(c/a) t}
=\frac{1}{a}(1-(b/a) \sqrt t-(c/a) t+(b/a)^2 t).
\]
Consequently, with this order of approximation, it follows for $j<k$ that
\[
\frac{1}{p_j}=\frac{1}{(r_j^0)^2}
\left[1- \frac{\Om_j}{r^0_j} \sqrt h  +\left(T+\frac{\Om_j^2-T_j}{(r^0_j)^2}\right)h \right].
\]
and the normalised states are
\[
\rho_j=\frac{\tilde \rho_j}{p_j}=[\rho -ih[H, \rho] -\frac12 h\{L^*L,\rho\})
+\frac{\sqrt h}{r^0_j} (\rho \tilde L_j^*+\tilde L_j \rho)+\frac{h}{(r^0_j)^2}\tilde L_j \rho \tilde L_j^*]
\]
\[
\times \left[1- \frac{\Om_j}{r^0_j} \sqrt h  +\left(T+\frac{\Om_j^2-T_j}{(r^0_j)^2}\right)h \right]
\]
\[
=\rho +  \frac{\sqrt h}{r^0_j}  (\rho \tilde L_j^* + \tilde L_j\rho-\Om_j \rho)+tB_j
\]
with
\[
B_j=-i[H, \rho] -\frac12 \{L^*L,\rho\}+ T \rho
 +(r_j^0)^{-2} (\tilde L_j\rho \tilde L_j^*-(\rho \tilde L_j^* + \tilde L_j\rho) \Om_j -T_j\rho +\Om_j^2  \rho).
\]

Therefore, the first $k$ terms in the expression for $(U_h-1)/h f(\rho)$, up to orders $h$ become
\[
\sum_{j=0}^{k-1} \frac{1}{h} ((r_j^0)^2 +\Om_j r^0_j \sqrt h)
[f(\rho +  \frac{\sqrt h}{r^0_j}  (\rho \tilde L_j^* + \tilde L_j\rho-\Om_j \rho)+hB_j)-f(\rho)]
\]
\[
=\sum_j  \frac{1}{h} ((r_j^0)^2 +\Om_j r^0_j \sqrt h)
f'(\rho) \left(\frac{\sqrt h}{r^0_j}  (\rho \tilde L_j^* + \tilde L_j\rho-\Om_j \rho)+hB_j\right)
\]
\[
+\frac12 \sum_j  \left(\rho \tilde L_j^* + \tilde L_j\rho-\Om_j \rho,
f''(\rho)  (\rho \tilde L_j^* + \tilde L_j\rho-\Om_j \rho)\right).
\]
The terms of order $1/\sqrt h$ cancel, because the coefficient at $h^{-1/2}$ equals
\[
\sum_{j=0}^{k-1} r_j^0   (\rho \tilde L_j^* + \tilde L_j\rho-\Om_j \rho)
\]
\[
=\sum_{l=1}^{k-1} \sum_{j=0}^{k-1} r^0_j r^l_j [\rho L_l +L_l \rho -{\tr}\, (\rho L_l+L_l \rho)]=0,
\]
due to the orthogonality of the vectors $r_j$. Consequently, disregarding terms of order $\sqrt h$, we have
\[
\sum_{j=0}^{k-1} \frac{1}{h} p_j (f(\rho_j)-f(\rho))
=\sum_j \left(f'(\rho), \Om_j  (\rho \tilde L_j^* + \tilde L_j\rho-\Om_j \rho)+(r^0_j)^2B_j\right)
\]
\[
+\frac12 \sum_j  \left(\rho \tilde L_j^* + \tilde L_j\rho-\Om_j \rho,
f''(\rho)  (\rho \tilde L_j^* + \tilde L_j\rho-\Om_j \rho)\right)
\]
\[
=\left(f'(\rho), -i[H, \rho] -\frac12 \{L^*L,\rho\} +\sum_{j=0}^{k-1}\tilde L_j \rho \tilde L_j^*
+\sum_{j=k}^{m-1} T_j \rho \right)
\]
\[
+\frac12 \sum_j  \left(\rho \tilde L_j^* + \tilde L_j\rho-\Om_j \rho,
f''(\rho)  (\rho \tilde L_j^* + \tilde L_j\rho-\Om_j \rho)\right).
\]
Noticing that
\[
L^*L=\sum_{j=1}^{k-1} L_j^*L_j=\tilde L^* \tilde L=\sum_{j=0}^{k-1}\tilde L_j^* \tilde L_j,
\]
we see that everything in the last expression rewrites in terms of $\tilde L$.

Including the terms with $j\ge k$ yields
\[
\lim_{h\to 0} \frac{U_h-1}{h} f(\rho)=\sum_{j=k}^{m-1} T_j \left(f(\frac{\tilde L_j \rho \tilde L_j^*}{T_j})-f(\rho)\right)
\]
\[
+\left(f'(\rho), -i[H, \rho] -\frac12 \{\tilde L^*\tilde L,\rho\}
+\sum_{j=0}^{k-1}\tilde L_j \rho \tilde L_j^*
+\sum_{j=k}^{m-1} T_j\right)
\]
\begin{equation}
\label{eqapprdif}
+\frac12 \sum_{j=0}^{k-1}  \left(\rho \tilde L_j^* + \tilde L_j\rho-\Om_j \rho,
f''(\rho)  (\rho \tilde L_j^* + \tilde L_j\rho-\Om_j \rho)\right),
\end{equation}
which is exactly the operator \eqref{eqmixgener} (for $\tilde L$).

Turning from $m-1$ operators $L_j$
to $m$ operators $\tilde L_j$ produces natural degeneracy meaning
that one can expect that often the number of terms in the last expression
can be made less than $k$. For instance, if $k=2$, $r_j=e_j$ for $j>1$ and $r_0,r_1$ belong to the
space generated by $e_0,e_1$, then $\tilde L_0=r_0^1 L_1$ and $\tilde L_1=r^1_1 L_1$
are proportional and
\[
\sum_{j=0}^{k-1} \tilde L_j \rho \tilde L_j^*=L_1\rho L_1^*,
\]
so that the corresponding diffusive part of \eqref{eqapprdif}
can be written as a single term (not as the sum of two terms):
\[
\sum_{j=0}^1 \frac{1}{h} p_j (f(\rho_j)-f(\rho))
=f'(\rho) \left(-i[H, \rho] -\frac12 \{\tilde L_1^*\tilde L_1,\rho\} +\tilde L_1 \rho \tilde L_1^* \right)
\]
\[
+\frac12 (\rho \tilde L_1^* + \tilde L_1\rho-{\tr} \, (\rho \tilde L_1^* +\tilde L_1 \rho) \rho) f''(\rho)
 (\rho \tilde L_1^* + \tilde L_1\rho-{\tr}\, (\rho \tilde L_1^* +\tilde L_1 \rho)  \rho).
\]

In \cite{KolQuantFrac} we suggested a different way of organising mixed observations via
explicit different channels, where the number of terms in the corresponding generator
is explicitly fixed from the initial model.

It is remarkable that the expression for the generator is continuous
in $r_j^0$ and has a finite limit as these coefficients
tend to zero, but this limit does not equal the generator
obtained in the case when all but one of these coefficients vanish.
The mere fact that they do not vanish creates the diffusive (second order) term in the generator.

Everything is ready for the main result of this Section:

\begin{theorem}
\label{thBeleqdif}

Let the operators $\tilde L_j$ be unitary for $j\ge k$. Then

(i) The scaled discrete semigroups $(U_h)^{[s/h]}$ converge to the semigroup $\Phi^{mix}_s$
from Theorem \ref{nonLinsdesemmix},
as $h\to 0$, so that the corresponding
processes converge in distribution,
 with the following rates of convergence:
\begin{equation}
\label{eq1thBeleqdif}
\|(U_h)^{[t/h]} -\Phi^{mix}_tf\| \le \sqrt h t M e^{Mt} \|f\|_{C^4(S(\HC))},
\end{equation}
with a constant $M$.

(ii) The scaled semigroups $T_t^h$ converge to the semigroup $\Phi^{mix}_t$,
as $h\to 0$, so that the corresponding
processes converge in distribution, with the following rates of convergence:
\begin{equation}
\label{eq2thBeleqdif}
\|T_t^hf -\Phi_t^{mix}f\| \le \sqrt h t M e^{Mt} \|f\|_{C^3(S(\HC))}.
\end{equation}
\end{theorem}

\begin{proof}
This is a consequence of \eqref{eqapprdif}, Theorem \ref{nonLinsdesemmix},
Proposition \ref{propconvsemigr}, and the observation that \eqref{eq5propconvsemigr}
holds here with  the triple of spaces $C^4(S(\HC))\subset C^2(S(\HC)) \subset C(S(\HC))$.
\end{proof}

\subsection{Appendix: convergence of semigroups}
\label{appsem}

Here we collect the results on the convergence of Markov semigroups and CTRWs, which form
the theoretical basis for our derivations of the filtering equations.

It is well known that the convergence of the generators on the core of the limiting generator
implies the convergence of semigroups. We shall use a version of this result with the rates,
namely the following result, given in Theorem 8.1.1 of \cite{Kolbook11}.

\begin{prop}
\label{propconvsemigr}

 Let $F_t=e^{tL}$ be a strongly continuous semigroup in a Banach space $B$ with a norm $\|.\|_B$,
 generate by an operator $L$,
having a core $D$, which is itself a Banach space with a norm $\|.\|_D\ge \|.\|_B$ so that $L\in \LC(D,B)$.
Let $F_t$ be also a bounded semigroup in $D$
such that $\|F_t\|_{D\to D} \le C_D(T)$ with a constant $C_D(T)$ uniformly for $t\in [0,T]$.

(i) Let $F_t^h$, $h>0$,  be a family of strongly continuous contraction semigroups in a Banach space $B$
with bounded generators $L_h$ such that
\[
\|L_hf-Lf\|_B \le \ep_h \|f\|_D
\]
for all $f\in D$ and some $\ep_h$ such that $\ep_h\to 0$ as $h\to 0$.
 Then the semigroups $F_t^h$ converge strongly to the semigroup $F_t$, as $h\to 0$, and
\begin{equation}
\label{eq1propconvsemigr}
\|F_t^hf -F_tf\|_B \le t \ep_h C_D(T)\|L\|_{D\to B}.
\end{equation}

(ii) Let $U_h$ be a family of contractions in $B$ such that
\begin{equation}
\label{eq2propconvsemigr}
\|\left(\frac{U_h-1}{h} -L\right)f\|_B \le \ep_h \|f\|_D,
\end{equation}
and
 \begin{equation}
\label{eq3propconvsemigr}
\|\left(\frac{F_h-1}{h} -L\right)f\|_B \le \ka_h \|f\|_D,
\end{equation}
with $\ep_h \to 0$ and $\ka_h\to 0$, as $h\to 0$.
Then the scaled discrete semigroups $(U_h)^{[t/h]}$ converge to the semigroup $F_t$
and moreover
 \begin{equation}
\label{eq4propconvsemigr}
\sup_{s\le t}\|(U_h)^{[s/h]} -F_sf\|_B \le (\ka_h+\ep_h)t \|f\|_B.
\end{equation}
\end{prop}

Additional condition \eqref{eq3propconvsemigr} makes working with discrete approximation a bit more subtle,
than with the continuous chain approximations. Effectively to get \eqref{eq3propconvsemigr}
one needs a deeper regularity. Namely one should have another core $\tilde D$ such that $D\subset \tilde D\subset B$
with $L\in \LC(D,\tilde D) \cap \LC(\tilde D,B)$. In this case it is easy to see that
\begin{equation}
\label{eq5propconvsemigr}
\|\left(\frac{F_h-1}{h} -L\right)f\|_B \le h \|L\|_{D,\tilde D} \|L\|_{\tilde D,B}\|f\|_D.
\end{equation}


\begin{thebibliography}{99}
\bibitem[Anti-Zeno]{Anti-Zeno} A.P.Balachandran, S.M.Roy A Quantum Anti-Zeno Paradox. Phys.Rev.Lett. 84 (2000), 4019-4022.
\bibitem[Attal]{Attal} S Attal and Y. Pautrat. From repeated to continuous quantum interactions. Ann. Henri Poincar\'e 7 (2006), 59--104.
\bibitem[BarbRock16]{BarbRock16} V. Barbu, M R\"ockner and D. Zhang. Stochastic nonlinear Schrödinger equations. Nonlinear Anal. 136 (2016), 168 - 194.
\bibitem[BarchBel]{BarchBel} A. Barchielli and V.P. Belavkin. Measurements contunuous in time and a posteriori states in quantum mechanics. J. Phys A: Math. Gen. 24 (1991), 1495-1514.
\bibitem[BarchBook]{BarchBook} A. Barchielli and M. Gregoratti. Quantum Trajectories and Measurements in Continuous Case. The Diffusive Case. Lecture Notes Physics, v. 782, Springer Verlag, Berlin, 2009.
\bibitem[BarchHol]{BarchHol} A. Barchielli and A.S. Holevo. Constructing quantum measurement processes via classical stochastic calculus. Stochastic Processes and their Applications 58 (1995) 293 - 317.
\bibitem[BarndLoub]{BarndLoub} O. E. Barndorff-Nielsen and E. R. Loubenets. General framework for the behaviour of continuously observed open quantum systems. J. Phys. A: Math. Gen. 35 (2002) 565 - 588.
\bibitem[Bel87]{Bel87} V. P. Belavkin, Nondemolition measurement and control in quantum dynamical systems. In: Information Complexity and Control in Quantum Physics. CISM Courses and Lectures 294, S. Diner and G. Lochak, eds., Springer-Verlag, Vienna, 1987, pp. 331–336.
\bibitem[Bel88]{Bel88} V.P. Belavkin. Nondemolition stochastic calculus in Fock space and nonlinear filtering and control in quantum systems. Proceedings XXIV Karpacz winter school (R. Guelerak and W. Karwowski, eds.), Stochastic methods in mathematics and physics. World Scientific, Singapore, 1988, pp. 310 - 324.
\bibitem[Bel92]{Bel92} V.P. Belavkin. Quantum stochastic calculus and quantum nonlinear filtering. J. Multivar. Anal. 42 (1992), 171 - 201.
\bibitem[BelKol]{BelKol} V.P. Belavkin, V.N. Kolokoltsov. Stochastic evolution as interaction representation of a boundary value problem for Dirac type equation. Infinite Dimensional Analysis, Quantum Probability and Related Fields 5:1 (2002), 61-92.
\bibitem[BoutHanJamQuantFilt]{BoutHanJamQuantFilt} L. Bouten, R. Van Handel and M. James. An introduction to quantum filtering. SIAM J. Control Optim. 46:6 (2007), 2199-2241.
\bibitem[Breuer]{Breuer} H.-P. Breuer and F. Petruccione. The theory of open quantum systems. Oxford University Press, 2002.
\bibitem[ChebFagn]{ChebFagn} A.M. Chebotarev and F. Fagnola. Sufficient conditions for conservativity of minimal quantum dynamical semigroups. J. Funct. Anal. 153 (1998), 382 - 404.
\bibitem[ChebQuez]{ChebQuez} A.M. Chebotarev, J. Garcia and R. Quezada. A priori estimates and existence theorems for the Lindblad equation with unbounded time-dependent coefficients. In Recent Trends in Infinite Dimensional Non-Commutative Analysis 1035 (1998), 44–65. Publ. Res. Inst. Math. Sci., Kokyuroku, Japan.
\bibitem[Eells]{Eells} J. Eells, A setting for global analysis. Bull. Amer. Math. Soc. 72 (1966), 751–807.
\bibitem[Fagnola]{Fagnola} F. Fagnola and C. M. Mora. Stochastic Schr\"odinger equation and applications to Ehrenfest-type theorems. ALEA Lat. Am. J. Probab. Math. Stat. 10:1 (2013), 191 - 223.
\bibitem[GreckGenerNonlinSchr]{GreckGenerNonlinSchr} W. Grecksch and H. Lisei. Stochastic nonlinear equations of Schr\"odinger type. Stoch. Anal. Appl. 29:4 (2011), 631 - 653.
\bibitem[Hayek]{Hayek} P. H\'ajek, V. Montesinos and V. Zizler. Geometry and Gˆateaux smoothness in separable Banach spaces. Operators and Matrices 6:2 (2012), 201–232.
\bibitem[Holevo91]{Holevo91} A.S. Holevo. Statistical Inference for quantum processes. In: Quanum Aspects of Optical communications. Springer LNP 378 (1991), 127-137, Berlin, Springer.
\bibitem[Holevo96]{Holevo96} A.S. Holevo. On dissipative stochastic equations in a Hilbert space. Probab. Theory Relat. Fields 104 (1996), 483 - 500.
\bibitem[Kallen]{Kallen} O. Kallenberg. Foundations of Modern Probability. Second ed., Springer, 2002.
\bibitem[Kol25a]{Kol25a} Vassili N. Kolokoltsov. On quantum stochastic master equations. Electron. J. Probab. 30 (2025), PNO = 57, 1-21.
\bibitem[Kol25b]{Kol25b} Vassili N. Kolokoltsov. On the Mathematical Theory of Quantum Stochastic Filtering Equations for Mixed States. Russian Journal of Mathematical Physics (RJMP) 32:3 (2025), 510-529.
\bibitem[KolQuantFrac]{KolQuantFrac} V. N. Kolokoltsov. Continuous time random walks modeling of quantum measurement and fractional equations of quantum stochastic filtering and control. Fractional Calculus and Applied Analysis 25 (2022), 128 - 165.
\bibitem[KolQuantLLN]{KolQuantLLN} V. N. Kolokoltsov. The law of large numbers for quantum stochastic filtering and control of many particle systems. Theoretical and Mathematical Physics 208:1 (2021), 97-121.
\bibitem[Kolbook11]{Kolbook11} V. N. Kolokoltsov. Markov processes, semigroups and generators. DeGruyter Studies in Mathematics v. 38, DeGruyter, 2011.
\bibitem[Kolbook19]{Kolbook19} V. N. Kolokoltsov. Differential equations on measures and functional spaces. Birkh\"auser Advanced Texts, Birkh\"auser, 2019.
\bibitem[Lakshm]{Lakshm} V. Lakshmikantham, R. Mutchell and R. W. Mitchell. Differential equations in closed subsets of a Banach space. Transactions of the American Mathematical Society 220 (1976), 103--113.
\bibitem[Loubenets]{Loubenets} E. R. Loubenets. Quantum Stochastic Approach to the Description of Quantum Measurements. J. Phys. A: Math. Gen. 34 (2001), 7639.
\bibitem[Martin]{Martin} R. H. Martin, Jr. Differential Equations on Closed Subsets of a Banach Space. \emph{Transactions of the American Mathematical Society.} \textbf{179} (1973), 399--414.
\bibitem[Mora13]{Mora13} C. M. Mora. Regularity of solutions to quantum master equations: a stochastic approach. The Annals of Probability 41:3B (2013), 1978 - 2012.
\bibitem[MoraRebo]{MoraRebo} C. M. Mora and R. Rebolledo. Basic Properties of Nonlinear Stochastic Schrödinger Equations Driven by Brownian Motions. The Annals of Applied Probability 18:2 (2008), 591 - 619.
\bibitem[Pellegrini]{Pellegrini} C. Pellegrini. Poisson and Diffusion Approximation of Stochastic Schr\"odinger Equations with Control. Ann. Henri Poincar\'e 10:5 (2009), 995–1025.
\bibitem[Pellegrini10]{Pellegrini10} C. Pellegrini. Markov chains approximation of jump–diffusion stochastic master equations. Ann. Henri Poincar\'e 46: 4 (2010), 924–948.
\bibitem[Phelps]{Phelps} R. R. Phelps. Convex Functions, Monotone Operators and Differentiability. Lecture Notes in Mathematics, v. 1364. Springer-Verlag Berlin Heidelberg 1993.
\bibitem[ReedSimon]{ReedSimon} M. Reed and B. Simon. Methods of modern mathematical physics. v .2. Academic Press, New York 1978.
\bibitem[Toolbox]{Toolbox} A. J. Guirao, V. Montesinos and V. Zizler. Renormings in Banach Spaces. A Toolbox. Springer Nature Switzerland AG, 2022.
\bibitem[VanNeerven]{VanNeerven} J. van Neerven, M. Veraar and L. Weis (2015). Stochastic Integration in Banach Spaces – a Survey. In: Dalang, R., Dozzi, M., Flandoli, F., Russo, F. (eds) Stochastic Analysis: A Series of Lectures. Progress in Probability, vol 68. Birkhäuser, Basel. %$https://doi.org/10.1007/978-3-0348-0909-2_11$
\bibitem[Yajima]{Yajima} K. Yajima. Existence of Solutions for Schrodinger Evolution Equations. Commun. Math. Phys. 110 (1987), 415-426.
\bibitem[Zeno]{Zeno} Misra, E. C. G. Sudarshan. The Zeno's paradox in quantum theory. J. Math. Phys. 18:4 (1977), 756–763.
\end{thebibliography}
\end{document}
